\documentclass[11pt]{article}
\usepackage[margin=1in]{geometry}
\usepackage[T1]{fontenc}
\usepackage[utf8]{inputenc}
\usepackage{lmodern}
\usepackage{amsmath,amssymb,amsthm,mathtools}
\usepackage{algorithm}
\usepackage{algpseudocode}
\usepackage{enumitem}
\usepackage{booktabs}
\usepackage[hidelinks]{hyperref}

\DeclareMathOperator*{\argmin}{arg\,min}
\newcommand{\E}{\mathbb{E}}
\newcommand{\Real}{\mathbb{R}}
\newcommand{\Reg}{\mathrm{Reg}}
\newcommand{\ind}[1]{\mathbf{1}_{\{#1\}}}
\theoremstyle{plain}
\newtheorem{Thm:Theorem}{Theorem}[section]
\newtheorem{Thm:Lemma}[Thm:Theorem]{Lemma}
\newtheorem{Thm:Proposition}[Thm:Theorem]{Proposition}
\newtheorem{Thm:Corollary}[Thm:Theorem]{Corollary}
\theoremstyle{remark}
\newtheorem{Thm:Remark}[Thm:Theorem]{Remark}
\numberwithin{equation}{section}
\setlength{\emergencystretch}{3em}

\title{Deadline-Constrained Online Allocation}
\author{Kamiar Asgari\thanks{AI usage disclaimer: Apart from the overarching idea for the problem setup, which was provided by the author, this paper was generated entirely by AI using Astra GPT-6 at the author's request.}}
\date{}

\begin{document}
\maketitle

\begin{abstract}
We study online allocation in a lossless scalar storage system with
adversarial arrivals in $[0,1]$ and a hard service deadline. Each arrival
must be consumed fully within the next $H$ actions. The controller chooses
its action before observing the current arrival and a jointly convex cost
of storage and consumption. Over a horizon of $T\ge H$ rounds, we measure regret
against the best fixed allocation of every arrival across the $H$
admissible service ages. Our controller organizes service ages in a balanced
binary tree and chooses each allocation split only when its interval
begins. This progressive allocation guarantees deadline feasibility and
regret $O(L_u\sqrt{HT}+L_xH^{3/2}\sqrt T)$, where
$L_x,L_u$ bound the state and action components of
joint subgradients. We also derive a prediction-error guarantee for causal
forecasts and show how uniformly positive storage marginal costs can
justify a shorter deadline, with an explicit correction for the finite
cost horizon. A lower bound matches the storage-dependent term when state
marginal costs may have either sign. For consumption-only costs, we prove
an $\Omega(L_u\sqrt T)$ lower bound; the optimal dependence on the deadline
remains open.
\end{abstract}

\section{Introduction}
\label{ch6:sec:introduction}

Online allocation with storage requires a controller to decide both how
much resource to use and when to use it. Resource feasibility alone does
not ensure timely service: a controller may consume only available
resource while deferring its use indefinitely. Even a bound on expected
backlog need not guarantee a maximum waiting time for each arrival. A hard
service deadline imposes an individual obligation: every arrival must be
fully served within a prescribed number of rounds. This paper studies how
to learn from changing costs while meeting that obligation.

We consider a lossless scalar storage system with arrivals in $[0,1]$.
Each arrival becomes available in the next round and must be consumed
completely within the next $H$ actions. At each round, the controller
chooses its consumption before observing the current arrival and cost
function. The cost is jointly convex in the stored resource and the
action, allowing holding and consumption costs to interact. Both arrivals
and costs may be adversarial.

The benchmark consists of policies that allocate the same fraction of
every arrival to each admissible service age. With a two-round deadline,
for example, a fixed policy could consume half of each arrival in the next
round and half in the following round. Regret is the difference between
the controller's cumulative cost and that of the best fixed age-allocation
policy chosen in hindsight. Each policy is evaluated on its own storage
and consumption trajectory under the same realized arrivals and costs.

Delay-constrained packet scheduling provides a concrete motivation.
In a fluid communication buffer, arrivals represent incoming data,
consumption represents transmission, and the deadline bounds the service
delay of each arrival. Spreading transmissions over time can reduce a
convex energy cost, but every arrival must still be transmitted on time.
Fixed age-allocation policies correspond to scheduling filters.
Section~\ref{ch6:subsec:delay-scheduling} develops this connection and its
modeling assumptions.

The central difficulty is that a service decision affects costs over
several rounds. Delaying consumption changes both its service cost and
the intervening storage costs. Committing to an entire schedule upon
arrival would fix every allocation before any of these costs is observed.
Our controller instead organizes the service ages in a balanced binary
tree. At each internal node, it divides the assigned resource between
earlier and later age intervals, choosing the split only when that node's
interval begins. Later splits can therefore use newly observed costs.
Conservation at every split guarantees complete service by the deadline,
while delayed mirror-descent updates learn how to allocate resource
between the two child intervals.

\subsection{Contributions and organization}

The analysis combines established scheduling-filter and sequential-learning
ideas with a causal service-age construction. In particular, the fixed
filter class has precedents in deadline scheduling
\cite{khojastepour2004delay}, and the local-to-global regret identity
specializes laminar regret decomposition \cite{farina2019laminar}.
Our claims concern the resulting guarantees for the storage model and
benchmark specified below; Section~\ref{ch6:sec:related} details the
closest overlaps and distinctions.

Let $T$ be the cost horizon, assume $1\le H\le T$, and let $L_x,L_u$ bound
the state and action components of joint subgradients. Our main results
are as follows.
\begin{enumerate}[leftmargin=*]
\item \textbf{Deadline feasibility and adversarial regret.}
The controller serves every arrival by its deadline and competes with the
full class of fixed age-allocation policies. For jointly convex
state--action costs, Theorem~\ref{ch6:thm:general} gives the regret bound
\[
 O\!\left(L_u\sqrt{HT}+L_xH^{3/2}\sqrt T\right).
\]
For fixed $H,L_x,L_u$, this yields $O(\sqrt T)$ regret.
\item \textbf{Shorter deadlines under positive storage costs.}
If the storage marginal cost has a known positive lower bound, accumulated
holding costs can outweigh the benefit of delaying consumption. We give an
explicit service deadline $K\le H$ that can shorten the allocation
horizon while retaining the original benchmark.
Theorem~\ref{ch6:thm:positive} bounds the resulting regret and accounts for
arrivals near the end of the cost horizon through an additional comparison
term of at most $L_uK$ when $K<H$.
\item \textbf{Prediction and correction.}
Causal forecasts guide the current allocation split, and the observed
costs subsequently update the learner. Theorem~\ref{ch6:thm:master} bounds
regret in terms of errors in local marginal-cost comparisons between
earlier and later service. Section~\ref{ch6:sec:prediction} constructs such
predictors from cost forecasts and known state-cost components.
\item \textbf{Lower bounds and the remaining gap.}
For $H\ge3$ and $T\ge4$, Theorem~\ref{ch6:thm:lower-state} proves an
$\Omega(L_xH^{3/2}\sqrt T)$ lower bound using linear storage costs whose
marginals may have either sign. Thus the storage-dependent term is
unavoidable in the worst case. For consumption-only costs with $H,T\ge2$,
Theorem~\ref{ch6:thm:lower-action} gives an $\Omega(L_u\sqrt T)$ lower bound,
leaving a gap of $\sqrt{H}$ from the upper bound. The optimal deadline
dependence in this case remains open.
\end{enumerate}

The paper proceeds from related work in Section~\ref{ch6:sec:related}
to the model and applications in Sections~\ref{ch6:sec:model}
and~\ref{ch6:sec:applications}. Sections~\ref{ch6:sec:algorithm}
and~\ref{ch6:sec:analysis} present the controller and its regret analysis.
The shorter-deadline and prediction refinements follow in
Sections~\ref{ch6:sec:positive} and~\ref{ch6:sec:prediction}.
Sections~\ref{ch6:sec:implementation} and~\ref{ch6:sec:lower} address
implementation and lower bounds; Sections~\ref{ch6:sec:next}
and~\ref{ch6:sec:conclusion} discuss future directions and conclude.

\section{Related Work}
\label{ch6:sec:related}

The relevant literature spans deadline scheduling, online control,
optimization with memory, and regret minimization on trees. Several
ingredients used here have established precedents: finite scheduling
filters, simplex policy parameters, local regret decompositions, delayed
updates, and optimistic mirror descent. The results of this paper concern
their use in a specific combination: lossless storage with adversarial
arrivals, complete service of every arrival within a common deadline,
jointly convex costs revealed after the action, and regret against the
entire class of fixed age distributions evaluated on their own
trajectories. We distinguish that combination from its individual
ingredients below.

\subsection{Deadline-feasible scheduling filters}
\label{ch6:subsec:related-scheduling}

Khojastepour and Sabharwal~\cite{khojastepour2004delay} provide a particularly
close antecedent. They study power-efficient communication subject to a
maximum delay for each packet and represent feasible schedulers using
finite, potentially time-varying linear filters. Their time-invariant
class has nonnegative coefficients summing to one. With the appropriate
shift for when arrivals become available, this is precisely the fixed
age-allocation class in Section~\ref{ch6:subsec:fixed-policies}. Under
i.i.d.\ arrivals and their long-run Gaussian-channel power criterion,
their optimal time-invariant scheduler is the uniform moving-average
filter, independently of the arrival distribution. Thus
neither normalized scheduling filters nor their deadline-feasibility
interpretation originates in the present paper.

Xu et al.~\cite{xu2010filter} develop a related finite-filter viewpoint for
server scheduling, studying workload smoothing and response-time
constraints through frequency-domain analysis. These applications show
that the filter interpretation extends beyond communication buffers.
Our benchmark is the best fixed filter on the realized sequence, whereas
the controller constructs a time-varying service allocation. Its learning
guarantee concerns arbitrary revealed-after-action convex costs, including
costs of the queue itself. The specific additional mechanism is to choose
each interval split when that interval begins, so that later choices for
the same arrival can use intervening observations. The older filter
representation alone does not supply the regret analysis for these
progressively chosen splits.

\subsection{Energy-efficient processor and packet scheduling}

Yao, Demers, and Shenker~\cite{yao1995scheduling} established the classical
speed-scaling model: jobs have release times, workloads, and deadlines,
and processor power is a convex function of speed. Bansal, Kimbrel, and
Pruhs~\cite{bansal2007speed} develop competitive algorithms and bounds for
energy and temperature management. The consumption-only specialization
of our model has the same basic incentive to spread divisible service
over time, and a common relative deadline is a special case of job
deadlines. These works compare energy with an unrestricted offline
schedule under a prescribed power function. Our comparator is restricted
to fixed age distributions, while the cost function may change
adversarially and is observed only after service. Consequently, the
competitive ratios and the additive regret bounds answer different
questions. Their competitive guarantees do not directly establish
additive regret to the fixed-filter benchmark under our information pattern.

For wireless transmission, Uysal-Biyikoglu, Prabhakar, and
El Gamal~\cite{uysal2002energy} analyze energy-efficient packet scheduling
and lazy transmission policies. Fu, Modiano, and
Tsitsiklis~\cite{fu2006transmission} study energy and deadline constraints
over fading channels, including a model in which current channel
information is unavailable until after transmission. Lee and
Jindal~\cite{lee2009energy} study the complete transmission of a fixed
amount of data within a finite number of slots under causal channel
information. The finite-horizon stochastic formulations of Fu et al.\
and Lee and Jindal establish important precedents for allocating service
across a deadline window. Our formulation has overlapping arrival cohorts
and adversarial revealed costs, with complete service imposed
deterministically. Post-action revelation by itself is therefore not
the distinction from all earlier packet-scheduling work.

\subsection{Learning and prediction under service deadlines}

Learning-based delay-constrained scheduling also predates this work.
Salodkar et al.~\cite{salodkar2008learning} learn energy-efficient policies
over fading channels without knowing the arrival and channel
distributions, using stochastic approximation with an average-delay
constraint. An average-delay guarantee does not imply the individual
completion requirement used here. Cayci and
Eryilmaz~\cite{cayci2019coding} analyze learning for dynamic coding with
deadline-constrained traffic and unknown channel statistics. Their regret
concerns the revenue of channel-allocation policies, including the
possibility of untransmitted packets, rather than obligatory complete
service of each arrival.

More directly, Mlika, Driouch, and Ajib~\cite{mlika2019deadline} combine
competitive and learning algorithms for deadline scheduling in
energy-harvesting networks. Their bandit learner selects among a finite
set of frame-level energy-allocation rules to increase the number of users
served on time. This is a precedent for regret guarantees in deadline
scheduling, but users may be dropped and the benchmark is a finite
collection of rules. Liu and Fang~\cite{liu2023scheduling} learn which task
types to activate when completion times and rewards are unknown; their
deadline mechanism interrupts tasks, together with a throughput
constraint. Neither interpretation of a deadline is equivalent to
conservation of every divisible arrival across its allowed service ages.

Predictions have also been studied in models with full service and common
relative deadlines. Bamas et al.~\cite{bamas2020learning} give
learning-augmented speed-scaling algorithms, including a model in which
each workload must be processed within a common delay. Their guarantees
interpolate between accurate workload predictions and worst-case
competitive performance. Antoniadis, Jabbarzade, and
Shahkarami~\cite{antoniadis2022prediction} develop a different prediction
setup for online speed scaling. These works optimize a known energy law
and assess error in predicted jobs or workloads. Here, forecasts predict
local comparisons of marginal service and storage costs, and the bound
is additive regret against fixed age filters. Shortening a virtual
deadline to create scheduling slack, as in learning-augmented speed
scaling, is also distinct from our comparison argument that bounds the
cost of shortening the comparator's support using positive holding
marginals.

Kong et al.~\cite{kong2025deadline} study deadline-aware fine-tuning on
spot and on-demand computing instances, with forecasts of market
conditions and regret for selecting among scheduling policies. The
finite policy-selection benchmark and computing-allocation model differ
from our continuum of age distributions. In another recent preprint,
Zhang and Mo~\cite{zhang2026hard} combine per-slot safety for a fluid
peak-age-of-information state with static regret under a viability
assumption. Their modeled age safety requires additional assumptions to
imply packet-level guarantees. It does not express the same obligation
as completing every arrival in the present storage system. The
distinction lies in the service obligations, information pattern, and
comparator, rather than in the general combination of deadlines and
no-regret learning.

\subsection{Deferrable loads, queueing, and energy storage}

Neely, Saber Tehrani, and Dimakis~\cite{neely2010renewable} allocate
renewable energy to delay-tolerant consumers using Lyapunov optimization
and a virtual queue that yields a worst-case delay bound. Their model
allows purchases from an outside supply and relates average performance
to the resulting delay bound, with extensions to arbitrary sample paths.
Thus queue-based control can guarantee maximum delay, not only average
backlog. Here the deadline is an exogenous requirement, arrivals are the
resource to be served, and the performance comparison is finite-horizon
regret to a fixed scheduling filter.

Gan et al.~\cite{gan2013deferrable} use model predictive control for
deferrable loads with arrival uncertainty and updated renewable-generation
forecasts. Their model allows heterogeneous deadlines and charging-rate
limits and studies aggregate-load variability. Lechowicz
et al.~\cite{lechowicz2026smoothed} study online smoothed demand management
with energy storage, flexible demand deadlines, and penalties for changing
purchase and delivery rates. They obtain competitive guarantees and
study learning from historical instances. These models have physical
features absent here, including additional capacity or smoothness
constraints. Our common-deadline, divisible-service model isolates a
different question: learning with arbitrary jointly convex state--action
costs while retaining exact per-arrival conservation.

Asgari and Neely~\cite{10.1145/3428337} combine Bregman-based online
optimization with energy-harvesting constraints and unknown i.i.d.\
energy arrivals. Their guarantees connect regret to battery capacity.
Chau, Zhang, and Chen~\cite{chau2016cost} study energy-storage management
through worst-case competitive guarantees against a fully informed
offline policy. Neither battery availability alone nor a bound on its
capacity enforces a maximum residence time for each arriving unit.
Conversely, compulsory consumption is not an innocuous assumption for
all energy-storage applications: the application must have a reason to
serve every arrival by its deadline.

\subsection{Simplex disturbance-action policies and nonstochastic control}
\label{ch6:subsec:related-control}

Agarwal et al.~\cite{agarwal2019online} develop online control with
adversarial disturbances and convex state--action costs using
disturbance-action policies, with regret against strongly stable linear
controllers. Golowich et al.~\cite{golowich2024population} use
simplex-structured policies for online population control. These works
establish both the disturbance-filter approach and the use of simplex
geometry in online control. In our problem the filter coordinates index
service ages and must conserve each arrival, rather than population
proportions or unconstrained disturbance responses.

The recent work of Asgari and Neely~\cite{asgari2026simplex} is especially
close: it treats scalar storage, adversarial nonnegative arrivals,
revealed-after-action joint costs, and fixed simplex disturbance-action
comparators. Its main dynamics have retention $\alpha\in(0,1)$; filter
coefficients form a subprobability simplex, and clipping ensures online
resource availability. Our fixed policies agree algebraically with the
unit-sum, $\alpha=1$ boundary of that filter family. However, its main
regret theorem relies on contraction and does not specialize to
$\alpha=1$. Its discussion of noncontracting storage instead adds a
stabilizing consumption baseline. Such a baseline excludes, for example,
for $H>1$, the pure age-$H$ policy before its service time. It also does not enforce
complete consumption of every arrival at a finite age. The present
deadline guarantee and comparison with all of $\Delta_H$ therefore
require a separate argument; the simplex representation and general
adversarial storage-learning formulation are shared antecedents.

\subsection{Exact constraints versus cumulative violation}

Online convex optimization with long-term constraints, exemplified by
Mahdavi, Jin, and Yang~\cite{mahdavi2012constraints}, trades regret against
cumulative constraint violation. Liu, Yang, and
Ying~\cite{liu2023constraints} study nonstochastic control with adversarial
and static constraints through constrained optimization with memory.
Bounds on cumulative violation cannot by themselves certify every
individual service deadline.

There are also online-control algorithms that maintain exact safety.
Li, Das, and Li~\cite{li2021affine} use tightened policy sets to satisfy
affine state and action constraints, under assumptions including a
strictly safe controller. Zhou and Tzoumas~\cite{zhou2023safe} study
nonstochastic control under time-varying convex safety constraints.
Thus exact feasibility with regret is already an established goal and
result in constrained control. Our distinction is the particular
conservation construction and comparator: every split preserves assigned
resource, leaf service exhausts it by the deadline, and all age
distributions remain admissible. The construction works from an empty
store and allows operation on the boundary, without a strict feasibility
margin. Merely expressing availability as an affine constraint does not
establish the required age-by-age completion property.

\subsection{Optimization with memory and the meaning of regret}
\label{ch6:subsec:related-memory}

Anava, Hazan, and Mannor~\cite{anava2015memory} study convex losses that
depend on several past decisions. Kumar, Dean, and
Kleinberg~\cite{kumar2023unbounded} extend this viewpoint to unbounded
memory and characterize regret through the influence of past decisions.
Their finite-memory specialization is also relevant to the deadline
dependence: Appendix D, Theorem D.1 gives a lower bound of order
$L h^{3/2}\sqrt T$ for a general memory-$h$ loss class with an
$\ell_1$ history norm. Accordingly, the exponent $3/2$ is not new as a
general memory phenomenon. The lower bound in
Theorem~\ref{ch6:thm:lower-state} realizes this order within a more
restricted class: scalar storage, linear signed state costs, exact
deadlines, and fixed age-policy comparators. A generic memory lower bound
does not automatically apply to every physically constrained subclass.

Memory methods also clarify why convexity of the fixed-policy objective
is insufficient for an immediate online algorithm. Updating a filter
parameter every round need not conserve the resource assigned by earlier
parameters; actual service and state are determined by past commitments.
The per-arrival representation explicitly records those commitments.
Zhao et al.~\cite{zhao2023nonstationary} study changing comparators through
dynamic policy regret in optimization with memory and nonstochastic
control. Our age-distribution comparator stays fixed; adversarially
changing costs and progressively changing online decisions do not turn
the guarantee into dynamic regret.

There is a further distinction when the environment reacts to actions.
Arora, Dekel, and Tewari~\cite{arora2012policy} distinguish ordinary regret
on the realized loss sequence from policy regret that reevaluates an
adaptive adversary along alternative histories. Here, the realized
arrival and cost sequence is held fixed for all comparisons, while each
fixed age policy generates its own state. Thus our guarantee addresses
the physical trajectory dependence for a fixed realized environment;
it does not compare different reactions of an adaptive environment to
counterfactual policies.

\subsection{Tree decompositions and local regret minimization}
\label{ch6:subsec:related-tree}

Zinkevich et al.~\cite{zinkevich2007counterfactual} introduced counterfactual regret
minimization for extensive-form games. Farina, Kroer, and
Sandholm~\cite{farina2019laminar} develop laminar regret decomposition
for sequential decision processes and convex losses. The principle is
to combine local learners at decision nodes, evaluating each alternative
with the learner's continuation decisions. Our binary tree identity is
a specialization of this local-to-global viewpoint: the comparator's
mass $M_I$ weights the local difference, and the node feedback compares
the two child continuation values. Computing hypothetical feedback even
at unused nodes is consistent with that interpretation. The identity
and the general idea of decomposing regret on a tree are not separate
novelty claims here.

Tree-based geometry and prediction have likewise been studied together.
Hoda et al.~\cite{hoda2010smoothing} construct smoothing geometries for
sequential-game strategy spaces. Farina et al.~\cite{farina2019stable}
combine predictive local learners with counterfactual regret
minimization, while Farina, Kroer, and
Sandholm~\cite{farina2019dilated,farina2025regularization} study optimistic
updates and improved regularization on sequential decision spaces.
Fan, Kroer, and Farina~\cite{fan2024optimality} analyze dilated entropy
and regret lower bounds for extensive-form games. These results provide context for
the local quadratic and entropy choices, but the associated strategy-space
or equilibrium guarantees do not directly give a physical service-age
bound. In particular, node-dependent feedback copies mean that our
implementation should not be identified without further argument with
one synchronous global mirror-descent step.

The additional structure is temporal. Nodes correspond to consecutive
service ages, their decisions begin at different physical times, and
feedback is complete only after their intervals end. The same interval
length $d_I$ controls both feedback delay and the storage-cost difference
between descendant ages. Balancing the tree makes the sums of
$\sqrt{d_I}$ and $d_I^{3/2}$ along a path scale respectively as
$\sqrt K$ and $K^{3/2}$ for a tree on $K$ ages. This is the source of the separate action and
storage contributions, rather than a generic dependence on the number
of tree nodes. The proof also has to link these local losses to the
actual jointly convex storage cost and to physically feasible service.

\subsection{Delayed feedback, optimism, and predictions}

Mirror descent and its Bregman geometry are classical
\cite{BECK2003167}. Joulani, Gy\"orgy, and
Szepesv\'ari~\cite{joulani2013delayed} provide general reductions for
learning with delayed feedback. The separate learner copies used here
are a bounded-delay device: a copy is reused only after its preceding
comparison has become available. The physical round cost is observed
immediately after service, but the complete comparison of two service
intervals depends on later rounds. It is this derived feedback, rather
than the original cost observation, that is delayed.
Quanrud and Khashabi~\cite{quanrud2015delays} and Joulani,
Gy\"orgy, and Szepesv\'ari~\cite{joulani2016delay} analyze other
delay-tolerant gradient and mirror-descent methods. Separate copies are
therefore a convenient implementation here, rather than a necessary
feature of all delayed-learning algorithms.

Rakhlin and Sridharan~\cite{RakhlinSridharan2013} develop learning with
predictable sequences. Flaspohler et al.~\cite{flaspohler2021optimism}
study optimism and delay together, showing how predictions can offset
missing feedback. These are direct methodological antecedents of our
prediction--correction bound. More recently, Abdullah
et al.~\cite{abdullah2026memory} combine constraints, memory, and
predictions through a delayed optimistic formulation, with regret and
cumulative-violation guarantees. The present contribution in this
direction is the choice of local marginal-cost prediction targets and
their aggregation under exact deadline feasibility. Predictions are
causal and may be inaccurate; the bound measures their realized local
errors rather than assuming access to future costs.
Cheng et al.~\cite{cheng2019predictor} also provide a predictor--corrector
approach to policy optimization with error-dependent regret. The
prediction--correction principle is shared; the service obligations and
local continuation comparisons are specific to the present construction.

\subsection{Inventory learning and budgeted online allocation}

Inventory learning provides another connection between optimization and
state-dependent implementability. Zhang, Chao, and
Shi~\cite{zhang2018perishable} study learning a base-stock policy for
perishable inventory under censored demand, with holding, lost-sales,
and outdating costs. Hihat et al.~\cite{hihat2023inventory} extend
inventory learning beyond i.i.d.\ demand under non-degeneracy conditions.
Pineci and Xu~\cite{pineci2026inventory} develop a hidden-target approach
for feasible inventory learning on general convex sets. In these
problems the controller replenishes stock to meet demand, and unsold
items may persist or perish. Here arrivals are exogenous and the
controller must consume all of them on time. A product's shelf life
therefore should not be identified with our mandatory service deadline.

Budgeted online allocation is also related in its use of mirror
descent, but the resource obligation is different. Balseiro, Lu, and
Mirrokni~\cite{balseiro2023allocation} study dual mirror descent for
reward-maximizing allocation under total resource budgets, with
guarantees across stochastic and adversarial input models. A total
budget limits aggregate resource expenditure; it does not require each
incoming unit to be allocated completely within a residence-time bound.
The shared term ``online allocation'' therefore covers different
feasibility structures and benchmarks.

\subsection{Scope of the contribution}

The contribution supported by these comparisons is a deadline-feasible
instantiation and analysis of progressive allocation for adversarial
joint costs, with the explicit bound
$O(L_u\sqrt{HT}+L_xH^{3/2}\sqrt T)$ against all fixed age filters.
The accompanying results identify a matching signed-storage lower bound
inside this model, construct local cost predictors, and bound the
finite-horizon cost of reducing the comparator's deadline under positive
holding marginals. We make no separate priority claim for scheduling
filters, simplex control, hierarchical regret decomposition, delayed
mirror descent, or the general coexistence of deadlines and learning.
The matching lower bound applies to signed state marginals; it does not
settle the optimal deadline dependence for consumption-only costs or for
the class with nonnegative state marginals.

\section{Model, Information, and Benchmark}
\label{ch6:sec:model}

We write $\Real_+$ for the nonnegative reals, $[n]=\{1,\ldots,n\}$,
$\E$ for expectation, and $\ind{A}$ for the indicator of an event $A$.
Bold symbols denote vectors, and $\partial f$ denotes the subdifferential.
All logarithms are natural unless a base is shown. Empty sums are zero,
and $0\log0=0$ in entropy expressions. For scalar clipping, write
$[v]_\ell^r=\min\{r,\max\{\ell,v\}\}$.

\subsection{Lossless dynamics and joint costs}
Let $T$ be the cost horizon and $H$ the consumption deadline, with integers
$1\le H\le T$ throughout.
Write $x_t,u_t,w_t$ for the stored resource, consumption, and arrival in
round $t$. The system starts empty and evolves as
\begin{equation}\label{ch6:eq:dynamics}
 x_1=0,\qquad x_{t+1}=x_t-u_t+w_t,\qquad
 0\le u_t\le x_t,\qquad w_t\in[0,1].
\end{equation}
At the beginning of round $t$, the controller knows the current state and
all arrivals and costs revealed through round $t-1$. It selects $u_t$ before
observing $w_t$ and $c_t$. After the action, the arrival and the entire cost
function are revealed, and the controller incurs $c_t(x_t,u_t)$.
The state in this cost is the beginning-of-round state; the arrival $w_t$
first enters $x_{t+1}$.

Each $c_t$ is a finite jointly convex function on
\begin{equation}\label{ch6:eq:domain}
 \mathcal{D}_H=\{(x,u):0\le u\le x\le H\}.
\end{equation}
Assume a joint-subgradient oracle: at every $(x,u)\in\mathcal{D}_H$, it can return
one pair $(\zeta^x,\zeta^u)$ such that
\begin{equation}\label{ch6:eq:support}
 c_t(x',u')\ge c_t(x,u)+\zeta^x(x'-x)+\zeta^u(u'-u)
 \quad\text{for all }(x',u')\in\mathcal{D}_H,
\end{equation}
and
\begin{equation}\label{ch6:eq:bounded-grad}
 |\zeta^x|\le L_x,\qquad |\zeta^u|\le L_u,
\end{equation}
for known finite constants $L_x,L_u\ge0$.
At the realized pair write
\begin{equation}\label{ch6:eq:actual-grad}
 \boldsymbol{\zeta}_t\coloneqq(\zeta_t^x,\zeta_t^u)^\top
 \in\partial c_t(x_t,u_t)
\end{equation}
for an oracle selection satisfying these bounds; its components are the state and action marginal costs.
For nonsmooth costs the two numbers must belong to the \emph{same joint
subgradient}; unrelated coordinatewise subgradients need not satisfy
\eqref{ch6:eq:support}.
Applying \eqref{ch6:eq:support} in both directions between two feasible
pairs and using \eqref{ch6:eq:bounded-grad} gives
\begin{equation}\label{ch6:eq:lipschitz}
 |c_t(x,u)-c_t(x',u')|\le L_x|x-x'|+L_u|u-u'|.
\end{equation}

The entire cost $c_t(x,u)$ may vary adversarially and may couple $x$ and $u$.
The residual-inventory example in
Section~\ref{ch6:subsec:coupled-penalties} illustrates an admissible
nonseparable cost and its signed marginals.
The regret bound holds for every realized arrival and cost sequence,
including sequences chosen in response to the controller. Once this
sequence has been produced, we keep it fixed when evaluating each
comparison policy. Each policy generates its own storage and consumption
trajectory, but all policies face the same arrivals and cost functions.

\subsection{Per-arrival service deadlines}
Arrival $w_s$ becomes available for the actions at times $s+1,s+2,\ldots$.
Let $q_{s,i}$ be its fraction actually consumed at age $i$. Thus $t$ indexes
rounds, $s$ arrivals, and $i$ ages. Feasible deadline accounting requires
\begin{equation}\label{ch6:eq:deadline}
 q_{s,i}\ge0,\qquad \sum_{i=1}^Hq_{s,i}=1,\qquad
 u_t=\sum_{i=1}^H w_{t-i}q_{t-i,i},
\end{equation}
with $w_s=0$ for $s\le0$. There is no service of $w_s$ outside
$s+1,\ldots,s+H$. The complete vector $\mathbf q_s$ need not be chosen upon arrival:
a causal controller may choose its fractions progressively, provided the
service at each time is determined before that time's cost is revealed.
For $w_s=0$, any such allocation vector can be used.

\begin{Thm:Proposition}[State identity and exact feasibility]\label{ch6:prop:state}
For any causal implementation of \eqref{ch6:eq:deadline},
\begin{equation}\label{ch6:eq:state}
 x_t=\sum_{i=1}^H w_{t-i}\sum_{j=i}^Hq_{t-i,j}.
\end{equation}
The dynamics \eqref{ch6:eq:dynamics} hold, every arrival is completely consumed
by its deadline, and $0\le u_t\le x_t\le H$. If, for some integer
$1\le K\le H$, all fractions have support in $1,\ldots,K$, then
$0\le u_t\le x_t\le K$.
\end{Thm:Proposition}
\begin{proof}
At the beginning of round $t$, the unserved part of arrival $w_{t-i}$ is
\[
w_{t-i}\left(1-\sum_{j=1}^{i-1}q_{t-i,j}\right)
=w_{t-i}\sum_{j=i}^{H}q_{t-i,j}.
\]
Summing over the possible ages gives \eqref{ch6:eq:state}, and
\begin{align*}
0\le u_t
&=\sum_{i=1}^{H}w_{t-i}q_{t-i,i}\\
&\le\sum_{i=1}^{H}w_{t-i}\sum_{j=i}^{H}q_{t-i,j}
=x_t.
\end{align*}
The state identity also gives
\begin{align*}
x_t-u_t+w_t
&=w_t+\sum_{i=1}^{H}w_{t-i}\sum_{j=i+1}^{H}q_{t-i,j}\\
&=w_t+\sum_{i=2}^{H}w_{t+1-i}\sum_{j=i}^{H}q_{t+1-i,j}\\
&=\sum_{i=1}^{H}w_{t+1-i}\sum_{j=i}^{H}q_{t+1-i,j}
=x_{t+1},
\end{align*}
using $\sum_{j=1}^{H}q_{t,j}=1$.  Moreover,
\[
0\le\sum_{j=i}^{H}q_{t-i,j}\le1,
\qquad
x_t\le\sum_{i=1}^{H}w_{t-i}\le H.
\]
If $q_{s,i}=0$ for $i>K$, the sums reduce to $1\le i\le K$, giving
$x_t\le K$.  Since $\sum_{i=1}^{H}w_sq_{s,i}=w_s$, every arrival is
exhausted by age $H$ (by age $K$ under the support restriction).
\end{proof}

\subsection{Fixed probability-simplex policies}
\label{ch6:subsec:fixed-policies}
Define
\begin{equation}\label{ch6:eq:simplex}
 \Delta_H=\left\{\mathbf m\in\Real_+^H:\sum_{i=1}^Hm_i=1\right\}.
\end{equation}
A fixed $\mathbf m\in\Delta_H$ assigns the same schedule to every arrival:
\begin{equation}\label{ch6:eq:fixed-pair}
 u_t(\mathbf m)=\sum_{i=1}^Hm_iw_{t-i},\qquad
 x_t(\mathbf m)=\sum_{i=1}^Hw_{t-i}\sum_{j=i}^Hm_j.
\end{equation}
Both coordinates are affine in $\mathbf m$, and the trajectory is feasible from
time one. The equivalent cumulative coordinates satisfy
$y_i=\sum_{j=1}^im_j$ and $0=y_0\le y_1\le\cdots\le y_H=1$.
In these coordinates, the state formula becomes
\[
 x_t(\mathbf m)=\sum_{i=1}^Hw_{t-i}(1-y_{i-1}),
\]
which expresses the stored resource through cumulative consumption,
with no storage decay and complete consumption by age $H$.

Write
\begin{equation}\label{ch6:eq:objective}
 J_T^{\mathrm{on}}=\sum_{t=1}^Tc_t(x_t,u_t),\qquad
 J_T(\mathbf m)=\sum_{t=1}^Tc_t(x_t(\mathbf m),u_t(\mathbf m)),
\end{equation}
and
\begin{equation}\label{ch6:eq:regret}
 \Reg_T(\Delta_H)=J_T^{\mathrm{on}}-\min_{\mathbf m\in\Delta_H}J_T(\mathbf m).
\end{equation}
Continuity from \eqref{ch6:eq:lipschitz} and compactness ensure existence of the
minimizer. The minimization chooses a common age distribution in hindsight.
All cost functions and arrivals
in \eqref{ch6:eq:objective} are the same realized functions and arrivals,
but each policy uses its own state.

\subsection{Service and costs after the horizon}\label{ch6:subsec:terminal}
Set $w_t=0$ for $t>T$ and execute all outstanding service through the actual
deadlines. Costs after $T$ are excluded from \eqref{ch6:eq:objective}; equivalently,
use zero cost functions and zero gradients there solely for analysis and
updates. If the costs of serving the remaining arrivals after the horizon are
part of the objective, include those rounds in $T$ for both policies.

For nontrivial learning let
\begin{equation}\label{ch6:eq:horizons}
 N=T-1.
\end{equation}
For each integer $1\le K\le H$, we identify a policy in $\Delta_K$ with its zero-padded
vector in $\Delta_H$ when evaluating its trajectory and $J_T$.
The arrival $w_T$ cannot affect any cost through $T$ and may be served at
$T+1$. Only arrivals $s=1,\ldots,N$ can affect costs included in the objective.
Every state within this horizon is also bounded by $\min\{H,t-1\}$.

For $T=1$ or $H=1$ regret is zero. The formulas below are stated for
$T\ge2$.

\section{Applications and Modeling Scope}
\label{ch6:sec:applications}

The model describes a divisible resource that can be stored without loss,
with every arrival served completely within a common deadline of $H$
rounds. There is no
independent service-capacity limit beyond resource availability in
\eqref{ch6:eq:dynamics}: several arrivals may be served together, and no
additional constraint $u_t\le1$ is imposed. Available storage must accommodate
the states permitted by \eqref{ch6:eq:domain}.

\subsection{Delay-constrained packet scheduling}
\label{ch6:subsec:delay-scheduling}

In a fluid communication buffer, $w_t$ is newly arriving data, $x_t$ is
the queued data before transmission, and $u_t$ is the amount successfully
transmitted in round $t$. An arrival at the end of round $s$ must be fully
transmitted during rounds $s+1,\ldots,s+H$. The constraint is therefore a
maximum delay for each arrival, in addition to the requirement that the
transmitter serve only data already in its buffer. Smoothing a burst over
several rounds can reduce a convex transmission-energy cost, but the delay
bound limits how long the scheduler can defer that service.

This is the scheduling setting motivating the connection to
Khojastepour and Sabharwal's \emph{Delay-Constrained Scheduling: Power
Efficiency, Filter Design, and Bounds}~\cite{khojastepour2004delay}.
The action of a fixed policy in \eqref{ch6:eq:fixed-pair} is a scheduling filter,
with weights in \eqref{ch6:eq:simplex}.
Nonnegative coefficients preserve causality and resource availability,
and their unit sum ensures complete transmission within $H$ rounds.
The uniform choice $m_i=1/H$ gives the moving-average schedule
$u_t=H^{-1}\sum_{i=1}^{H}w_{t-i}$.
The controller learns an adaptive allocation of each arrival across
these ages and measures its performance against the best fixed filter
on the realized sequence. Section~\ref{ch6:subsec:related-scheduling}
discusses the relationship to the earlier scheduling objective.

For a concrete energy-aware specialization, let
\[
 c_t(x,u)=\rho_t P(u)+\lambda x,
 \qquad 0\le\rho_t\le\bar\rho,\qquad \lambda\ge0,
\]
where $P$ is a known differentiable convex transmission-energy function
with $0\le P'(u)\le G_P$ on $[0,H]$. The coefficient $\rho_t$ is an
energy price or cost weight revealed after the action. The joint gradient
is $(\lambda,\rho_tP'(u))$, so one may take $L_x=\lambda$ and
$L_u=\bar\rho G_P$. The case $\lambda=0$ uses the consumption-only
guarantee in Corollary~\ref{ch6:cor:action}; Corollary~\ref{ch6:cor:constant}
applies to the constant nonnegative state marginal cost $\lambda$. Setting $\rho_t=1$ and
$\lambda=0$ recovers a fixed convex transmit-energy objective.

This specialization assumes successful service. Unknown channel capacity
and failed transmissions would change both the service feasibility and
the feedback.

\subsection{Coupled penalties and the role of signed marginals}
\label{ch6:subsec:coupled-penalties}

For a target residual inventory $d_t\in[0,H]$ revealed after service, consider
\[
 c_t(x,u)=\frac{(x-u-d_t)^2}{2H}.
\]
This jointly convex cost penalizes deviations of the residual inventory
$x-u$ from the target $d_t$. The residual is measured after service and
before the new arrival. The cost's state and action marginals have magnitudes at
most one, so $L_x=L_u=1$ is valid. The state marginal can be negative when
the residual is below the target. The general theorem applies, whereas the
nonnegative state-marginal refinement need not. This example shows why a nonnegative
\emph{cost value} is different from a nonnegative \emph{state marginal}.

The regret guarantee for all these applications uses the benchmark in
Section~\ref{ch6:subsec:fixed-policies}.

\section{Causal Allocation and Delayed Mirror Descent}
\label{ch6:sec:algorithm}
Choose a deadline $1\le K\le H$ for the controller. The general controller uses
$K=H$; Section~\ref{ch6:sec:positive} gives a justified reduction when
state marginals have a uniform positive lower bound.

\subsection{Age intervals and causal routing}
Build a balanced binary tree $\mathcal{T}_K$ on the ordered leaves $1,\ldots,K$.
Tree intervals $[a,b]$ denote the integer age sets $\{a,\ldots,b\}$.
For an internal node $I=[a,b]$, write $a_I=\min I=a$. Its length
and children are
\begin{equation}\label{ch6:eq:tree-def}
 d_I=b-a+1,\qquad r_I=\lfloor(a+b)/2\rfloor+1,\qquad
 I_{\mathrm{L}}=[a,r_I-1],\quad I_{\mathrm{R}}=[r_I,b].
\end{equation}
Let $\mathcal{I}_K$ be its set of internal nodes.
For arrival $s$, select $p_{s,I}\in[0,1]$ at the end of round $s+a-1$.
Route fraction $p_{s,I}$ of the resource assigned to $I$ to its left child,
and the remainder to its right child. Retain right-child resource until
its interval begins. Each leaf consumes all incoming resource at its age.
Simultaneously starting decisions are processed from ancestors to descendants.
The root receives $w_s$ when it arrives.

The fractions are deterministic divisions of divisible resource, not
randomized service choices. The final fraction
$q_{s,i}$ is the product of the routing factors on the path to leaf $i$.
All factors on that path have been chosen by the
end of $s+i-1$. They are nonnegative and their leaf products sum to one.
Proposition~\ref{ch6:prop:state} therefore proves causality, feasibility, and
complete service by age $K\le H$, independently of how the splits are learned.

\subsection{A four-age example}
\label{ch6:subsec:example}

Figure~\ref{ch6:fig:routing} illustrates the timing for $K=4$.
Suppose one unit arrives at the end of round $s$. At that time, the root
assigns half to ages $1,2$ and half to ages $3,4$. The left child
immediately assigns one quarter of its incoming resource to age one. The
right child's split can wait until the end of round $s+2$, after two
more costs have been observed. If that later split assigns three quarters
to age three, the final allocation vector is
\[
 \mathbf q_s=(q_{s,1},q_{s,2},q_{s,3},q_{s,4})^\top
 =\left(\frac18,\frac38,\frac38,\frac18\right)^\top.
\]
This example uses illustrative split fractions to show how a feasible
allocation is assembled as successive age intervals begin.

\begin{figure}[tb]
\centering
\setlength{\unitlength}{1mm}
\begin{picture}(138,67)
\put(44,59){\makebox(50,5){$[1,4]$: decide at end of $s$}}
\put(63,56){\vector(-3,-2){27}}
\put(75,56){\vector(3,-2){27}}
\put(37,49){\makebox(12,5){$1/2$}}
\put(88,49){\makebox(12,5){$1/2$}}
\put(3,30){\makebox(55,6){$[1,2]$: end of $s$}}
\put(78,30){\makebox(58,6){$[3,4]$: end of $s+2$}}
\put(27,29){\vector(-1,-1){15}}
\put(34,29){\vector(1,-1){15}}
\put(103,29){\vector(-1,-1){15}}
\put(111,29){\vector(1,-1){15}}
\put(9,22){\makebox(12,5){$1/4$}}
\put(41,22){\makebox(12,5){$3/4$}}
\put(85,22){\makebox(12,5){$3/4$}}
\put(117,22){\makebox(12,5){$1/4$}}
\put(0,7){\makebox(25,5){Age $1$}}
\put(37,7){\makebox(25,5){Age $2$}}
\put(76,7){\makebox(25,5){Age $3$}}
\put(113,7){\makebox(25,5){Age $4$}}
\put(0,0){\makebox(25,5){$q_{s,1}=1/8$}}
\put(37,0){\makebox(25,5){$q_{s,2}=3/8$}}
\put(76,0){\makebox(25,5){$q_{s,3}=3/8$}}
\put(113,0){\makebox(25,5){$q_{s,4}=1/8$}}
\end{picture}
\caption[Allocation over four consumption ages]{Allocation over four consumption ages. Edge labels
are fractions of the resource entering their parent. Leaf $i$ serves
its assigned resource at time $s+i$; the right subtree can use information
revealed after the earlier service decisions.}
\label{ch6:fig:routing}
\end{figure}

\subsection{Marginal feedback from the actual joint cost}
After each round within the cost horizon, obtain the joint pair
$(\zeta_t^x,\zeta_t^u)$ in \eqref{ch6:eq:actual-grad}; set both components to zero
after $T$. For arrival $s\le N$ define
\begin{equation}\label{ch6:eq:leaf}
 A_{s,i}=w_s\left(\zeta_{s+i}^u+\sum_{k=1}^{i}\zeta_{s+k}^x\right),
 \qquad i=1,\ldots,K.
\end{equation}
In the linearized total cost, $A_{s,i}$ is the coefficient of the fraction
of arrival $s$ served at age $i$. Its service term comes from round $s+i$,
while its storage terms come from rounds $s+1,\ldots,s+i$. The storage sum
includes the service round because the cost uses the beginning-of-round
state. The factor $w_s$ converts an allocation fraction into an amount of
resource. All subgradients are evaluated along the controller's realized
trajectory.

Define subtree values recursively by
\begin{equation}\label{ch6:eq:values}
 V_{s,\{i\}}=A_{s,i},\qquad
 V_{s,I}=p_{s,I}V_{s,I_{\mathrm{L}}}+(1-p_{s,I})V_{s,I_{\mathrm{R}}}.
\end{equation}
The value $V_{s,I}$ is an average of the marginal-cost coefficients for
ages in $I$, with weights determined by the controller's splits within that
node's subtree. It gives the linearized cost coefficient for routing an amount $w_s$
through $I$ according to those splits.
The scalar feedback is
\begin{equation}\label{ch6:eq:g}
 g_{s,I}=V_{s,I_{\mathrm{L}}}-V_{s,I_{\mathrm{R}}}.
\end{equation}
This feedback compares routing an amount $w_s$ through the earlier child
with routing that amount through the later child, using the observed cost
slopes and the controller's splits within each child's subtree. A positive value
favors sending less resource to the earlier interval on future uses of
this learner; a negative value favors sending more.

At each node, we compare the costs of routing the full arrival amount
through the earlier and later children. This is a hypothetical comparison,
so we calculate it even when the controller sends no resource to that
node. Split fractions must therefore be chosen in unused branches as
well. In the regret proof, each comparison is weighted by the fraction
that the fixed comparison policy assigns to that node, denoted $M_I$;
see Lemma~\ref{ch6:lem:tree}.

At a node $I=[a,b]$ spanning $d_I$ ages, complete feedback becomes
available $d_I$ rounds after the split is chosen. We therefore maintain
$d_I$ separate copies of the learning rule and use them in turn for
successive arrivals. Each copy receives feedback from its previous
decision before it is used again. In particular, the copy used for
arrival $s$ is next used for arrival $s+d_I$, at the end of round
\begin{equation}\label{ch6:eq:timing}
 (s+d_I)+a-1=s+b.
\end{equation}
Bounded-delay reductions are standard \cite{joulani2013delayed}.
The controller collects observations
during the interval and updates the learner when the full comparison is available.

\subsection{Updating a split with quadratic or entropy penalties}
Each learner stores a reference split $z\in[0,1]$, initially $1/2$.
To choose a split, it balances the predicted cost against a penalty for
moving away from $z$. The default penalty is the squared distance
$\tfrac12(p-z)^2$. We also analyze the binary relative entropy penalty
$\operatorname{kl}(p\|z)$. We write $D_\psi(p,z)$ when either penalty can
be used.
Explicitly,
\[
 D_{\psi_Q}(p,z)=\frac12(p-z)^2,\qquad
 D_{\psi_E}(p,z)=\operatorname{kl}(p\|z)
 =p\log\frac pz+(1-p)\log\frac{1-p}{1-z}.
\]
The entropy formula is defined for $p\in[0,1]$ and $z\in(0,1)$.
Appendix~\ref{app:geometry} gives the generating functions and their
basic bounds.
In this paper $\psi$ denotes the regularizer and $F_{s,I}$ the predictor.

The function $F_{s,I}(p)$ predicts how the local cost changes with the
fraction $p$ sent to the earlier interval. Its derivative at the chosen
split is the prediction of $g_{s,I}$ used in the regret bound. For example,
a forecast $\widehat g$ of this marginal comparison gives the linear
predictor $F_{s,I}(p)=\widehat g p$. The predictor must use only information
available when the split is chosen.

At the decision time choose a convex differentiable predictor
$F_{s,I}:[0,1]\to\Real$ with finite derivative satisfying this information
requirement, and choose
\begin{equation}\label{ch6:eq:play}
 p_{s,I}=\argmin_{p\in[0,1]}
 \{\eta_IF_{s,I}(p)+{D_\psi}(p,z_{s,I})\}.
\end{equation}
The basic algorithm uses $F_{s,I}=0$, so $p_{s,I}=z_{s,I}$.
When $g_{s,I}$ is complete, update the \emph{reference value stored by the same learner} by
\begin{equation}\label{ch6:eq:correct}
 z_{s+d_I,I}=\argmin_{z\in[0,1]}
 \{\eta_Ig_{s,I}z+{D_\psi}(z,z_{s,I})\}.
\end{equation}
Minimizing these scalar objectives gives the closed forms
(see Appendix~\ref{app:geometry})
\begin{align}
 z_{s+d_I,I}&=[z_{s,I}-\eta_Ig_{s,I}]_0^1,&&\text{quadratic},\label{ch6:eq:quad}\\
 z_{s+d_I,I}&=\frac{z_{s,I}\exp(-\eta_Ig_{s,I})}
 {1-z_{s,I}+z_{s,I}\exp(-\eta_Ig_{s,I})},&&\text{binary entropy}.\label{ch6:eq:entropy}
\end{align}
The prediction shifts the chosen split away from the reference $z_{s,I}$.
Once feedback arrives, the next reference incorporates the observed
marginal comparison starting from that original reference. For example,
with quadratic regularization, a linear predictor $F(p)=\widehat g p$,
and no active boundary constraint, the chosen split is
$p=z-\eta_I\widehat g$. Updating from $p$ would give
$z-\eta_I(\widehat g+g)$, retaining the predicted step in addition to the
observed step. Updating from $z$ gives $z-\eta_Ig$.

With the entropy penalty, the reference values and chosen splits remain strictly between zero and one for finite feedback and derivatives;
quadratic iterates may equal zero or one. Both the split-selection and reference-update minimizations have
unique solutions. The one-step bound is proved in
Lemma~\ref{app:lem:optimistic}.

\begin{algorithm}[tb]
\caption{Deadline allocation with joint costs}
\label{ch6:alg:controller}
\begin{algorithmic}[1]
\Require $T$, $K\le H$, tree $\mathcal{T}_K$, regularizer $\psi$, positive
 steps $\eta_I$, causal predictor rule $F$
\State If $K=1$, serve every arrival at age one and stop.
\State Initialize every node's $d_I$ reference copies to $1/2$;
 initialize resource and feedback records to zero.
\For{$t=1,\ldots,T+K$}
 \State Execute the due leaf services, whose sum is $u_t$.
 \If{$t\le T$}
  \State Observe $w_t,c_t$; incur $c_t(x_t,u_t)$; obtain one joint
  subgradient $(\zeta_t^x,\zeta_t^u)$ at $(x_t,u_t)$.
 \Else
  \State Set $w_t\gets0$, $\zeta_t^x\gets0$, and $\zeta_t^u\gets0$; add no new objective term.
 \EndIf
 \State Set $x_{t+1}\gets x_t-u_t+w_t$.
 \State Complete due leaf and subtree values using \eqref{ch6:eq:leaf}--\eqref{ch6:eq:values},
 descendants before ancestors.
 \State For each completed $(s,I)$, form $g_{s,I}$ and update its copy
 by \eqref{ch6:eq:correct}.
 \State For each starting $(s,I)$ with $1\le s\le N$ and
 $t=s+a_I-1$, choose the split in \eqref{ch6:eq:play} and route its mass,
 ancestors before descendants.
 \State If $t=T$, assign $w_T$ to service at $T+1$.
\EndFor
\end{algorithmic}
\end{algorithm}

\section{Regret Analysis for General Joint Costs}
\label{ch6:sec:analysis}

\subsection{Regularizer constants and a prediction-error bound at one node}
The local prediction--correction estimate uses standard optimistic
mirror-descent analysis \cite{RakhlinSridharan2013,flaspohler2021optimism}.
For the two regularizers define constants by
\begin{equation}\label{ch6:eq:geometry}
 {D_\psi}(p,z)\ge\frac{\sigma_\psi}{2}(p-z)^2,\qquad
 \sup_{p\in[0,1]}{D_\psi}(p,1/2)\le{A_\psi},\qquad
 {C_\psi}=\sqrt{2{A_\psi}/\sigma_\psi}.
\end{equation}
Valid values are
\begin{equation}\label{ch6:eq:constants}
 \begin{array}{c|ccc}
 &\sigma_\psi&{A_\psi}&{C_\psi}\\\hline
 \text{quadratic}&1&1/8&1/2\\
 \text{binary entropy}&4&\log2&\sqrt{\log2/2}
 \end{array}
\end{equation}
These constants follow from the scalar geometry bounds in
Appendix~\ref{app:geometry}.
Apply Lemma~\ref{app:lem:optimistic} with regularizer
$\psi$, predictor $F=F_{s,I}$, and feedback $g=g_{s,I}$.
Writing $\widehat g=F'_{s,I}(p)$ and dropping the nonpositive term
$-D_\psi(p,z)/\eta$ gives, for every $p^*\in[0,1]$,
\begin{equation}\label{ch6:eq:one-step}
 g(p-p^*)\le
 \frac{D_\psi(p^*,z)-D_\psi(p^*,z^+)}\eta
 +\frac\eta{2\sigma_\psi}(g-\widehat g)^2.
\end{equation}
The lemma covers quadratic boundary solutions; with the entropy penalty, the reference values and chosen splits
remain strictly between zero and one because feedback and predictor derivatives are finite.

\begin{Thm:Lemma}[Local regret with delayed feedback]\label{ch6:lem:local}
Set $e_{s,I}=g_{s,I}-F'_{s,I}(p_{s,I})$. For any fixed $p^*\in[0,1]$,
\begin{equation}\label{ch6:eq:local}
 \sum_{s=1}^Ng_{s,I}(p_{s,I}-p^*)
 \le\frac{d_I{D_\psi}(p^*,1/2)}{\eta_I}
 +\frac{\eta_I}{2\sigma_\psi}\sum_{s=1}^Ne_{s,I}^2.
\end{equation}
If a known $B_I>0$ bounds $|e_{s,I}|$ uniformly, the step size
\begin{equation}\label{ch6:eq:steps}
 \eta_I=\sqrt{\frac{2\sigma_\psi{A_\psi} d_I}{NB_I^2}}
\end{equation}
gives
\begin{equation}\label{ch6:eq:local-opt}
 \sum_{s=1}^Ng_{s,I}(p_{s,I}-p^*)
 \le{C_\psi} B_I\sqrt{d_IN}.
\end{equation}
\end{Thm:Lemma}
\begin{proof}
Fix $I$ and abbreviate $d=d_I$, $\eta=\eta_I$.  For
$r\in[d]$, let
\[
\mathcal S_r=\{s\in[N]:s\equiv r\pmod d\}
=\{s_{r,1}<\cdots<s_{r,n_r}\}.
\]
By \eqref{ch6:eq:timing},
$z_{s_{r,k}+d,I}=z_{s_{r,k+1},I}$ whenever $k<n_r$.
Summing \eqref{ch6:eq:one-step} along one nonempty copy therefore gives
\begin{align*}
\sum_{k=1}^{n_r}g_{s_{r,k},I}(p_{s_{r,k},I}-p^*)
&\le\frac{1}{\eta}
 \sum_{k=1}^{n_r}\left[
 D_\psi(p^*,z_{s_{r,k},I})
 -D_\psi(p^*,z_{s_{r,k}+d,I})\right]\\
&\quad+\frac{\eta}{2\sigma_\psi}
 \sum_{k=1}^{n_r}e_{s_{r,k},I}^2\\
&=\frac{D_\psi(p^*,1/2)
 -D_\psi(p^*,z_{s_{r,n_r}+d,I})}{\eta}
 +\frac{\eta}{2\sigma_\psi}
 \sum_{k=1}^{n_r}e_{s_{r,k},I}^2\\
&\le\frac{D_\psi(p^*,1/2)}{\eta}
 +\frac{\eta}{2\sigma_\psi}
 \sum_{k=1}^{n_r}e_{s_{r,k},I}^2.
\end{align*}
Summing over the at most $d$ nonempty residue classes gives
\eqref{ch6:eq:local}.
If $|e_{s,I}|\le B_I$, the step size \eqref{ch6:eq:steps} gives
\begin{align*}
\frac{d_IA_\psi}{\eta_I}
+\frac{\eta_I}{2\sigma_\psi}\sum_{s=1}^{N}e_{s,I}^2
&\le\frac{d_IA_\psi}{\eta_I}
+\frac{\eta_INB_I^2}{2\sigma_\psi}\\
&=2B_I\sqrt{\frac{A_\psi d_IN}{2\sigma_\psi}}\\
&=C_\psi B_I\sqrt{d_IN}.\qedhere
\end{align*}
\end{proof}
For any fixed step size, the regret bound depends on the prediction errors
that actually occur. These errors need not be known in advance. To choose
the step size before running the controller, we can use a known upper
bound $B_I$ on their magnitude, as in \eqref{ch6:eq:steps}.

\subsection{Joint convexity and the tree identity}
\begin{Thm:Lemma}[Regret linearization]\label{ch6:lem:joint}
For every $\mathbf m\in\Delta_K$,
\begin{equation}\label{ch6:eq:joint-compare}
 J_T^{\mathrm{on}}-J_T(\mathbf m)
 \le\sum_{s=1}^N\sum_{i=1}^KA_{s,i}(q_{s,i}-m_i).
\end{equation}
\end{Thm:Lemma}
\begin{proof}
Put $\delta_{s,i}=q_{s,i}-m_i$ and retain the convention
$\zeta_r^x=\zeta_r^u=0$ for $r>T$.  The subgradient inequality at the
realized state--action pair gives
\[
c_t(x_t,u_t)-c_t(x_t(\mathbf m),u_t(\mathbf m))
\le\zeta_t^x(x_t-x_t(\mathbf m))
+\zeta_t^u(u_t-u_t(\mathbf m)).
\]
The action contribution rearranges as
\begin{align*}
\sum_{t=1}^{T}\zeta_t^u\bigl(u_t-u_t(\mathbf m)\bigr)
&=\sum_{t=1}^{T}\sum_{i=1}^{K}
  \zeta_t^uw_{t-i}\delta_{t-i,i}\\
&=\sum_{s=1}^{N}\sum_{i=1}^{K}
  w_s\zeta_{s+i}^u\delta_{s,i}.
\end{align*}
Similarly, the state contribution is
\begin{align*}
\sum_{t=1}^{T}\zeta_t^x\bigl(x_t-x_t(\mathbf m)\bigr)
&=\sum_{t=1}^{T}\sum_{i=1}^{K}
  \zeta_t^xw_{t-i}\sum_{j=i}^{K}\delta_{t-i,j}\\
&=\sum_{s=1}^{N}w_s\sum_{i=1}^{K}
  \zeta_{s+i}^x\sum_{j=i}^{K}\delta_{s,j}\\
&=\sum_{s=1}^{N}\sum_{j=1}^{K}
  w_s\left(\sum_{i=1}^{j}\zeta_{s+i}^x\right)\delta_{s,j}.
\end{align*}
Adding these contributions and using \eqref{ch6:eq:leaf} proves the bound.
For fixed realized costs, each $(x_t,u_t)$ is affine in the final fractions;
therefore $\sum_t c_t(x_t,u_t)$ is jointly convex in those fractions.
\end{proof}

We can describe a fixed policy using the same tree as the controller.
At node $I$, $M_I$ is the fraction of each arrival assigned to the ages
in that interval. When $M_I>0$, $p_I^*$ is the proportion of that
fraction assigned to the earlier child:
\begin{equation}\label{ch6:eq:comparator-mass}
 M_I=\sum_{i\in I}m_i,\qquad
 p_I^*=\frac{M_{I_{\mathrm{L}}}}{M_I}\quad(M_I>0).
\end{equation}
If $M_I=0$, the split can be chosen arbitrarily because the policy assigns
no resource there.

The next identity expresses the linearized cost difference for one
arrival as a sum of comparisons at individual nodes. Each comparison is
weighted by the fraction $M_I$ assigned there by the fixed policy.
This is the binary terminal-cost specialization of the counterfactual
and laminar regret decomposition viewpoint
\cite{zinkevich2007counterfactual,farina2019laminar}; we include its
short proof to make the service-age analysis self-contained.

\begin{Thm:Lemma}[Tree identity]\label{ch6:lem:tree}
For every arrival $s$,
\begin{equation}\label{ch6:eq:tree-identity}
 \sum_{i=1}^KA_{s,i}(q_{s,i}-m_i)
 =\sum_{I\in\mathcal{I}_K}M_I(p_{s,I}-p_I^*)g_{s,I}.
\end{equation}
\end{Thm:Lemma}
\begin{proof}
If $M_I>0$, then $M_Ip_I^*=M_{I_{\mathrm{L}}}$ and
$M_I(1-p_I^*)=M_{I_{\mathrm{R}}}$.  Hence
\begin{align*}
M_I(p_{s,I}-p_I^*)g_{s,I}
&=(M_Ip_{s,I}-M_{I_{\mathrm{L}}})
  (V_{s,I_{\mathrm{L}}}-V_{s,I_{\mathrm{R}}})\\
&=M_I\!\left[p_{s,I}V_{s,I_{\mathrm{L}}}
 +(1-p_{s,I})V_{s,I_{\mathrm{R}}}\right]
  -M_{I_{\mathrm{L}}}V_{s,I_{\mathrm{L}}}-M_{I_{\mathrm{R}}}V_{s,I_{\mathrm{R}}}\\
&=M_IV_{s,I}-M_{I_{\mathrm{L}}}V_{s,I_{\mathrm{L}}}-M_{I_{\mathrm{R}}}V_{s,I_{\mathrm{R}}}.
\end{align*}
If $M_I=0$, then $M_{I_{\mathrm{L}}}=M_{I_{\mathrm{R}}}=0$, so the same identity holds for
any $p_I^*$.  Summing it over all internal nodes cancels every internal
child term:
\begin{align*}
\sum_{I\in\mathcal I_K}M_I(p_{s,I}-p_I^*)g_{s,I}
&=M_{\mathrm{root}}V_{s,\mathrm{root}}
  -\sum_{i=1}^{K}M_{\{i\}}V_{s,\{i\}}\\
&=V_{s,\mathrm{root}}-\sum_{i=1}^{K}m_iA_{s,i}.
\end{align*}
The routing recursion \eqref{ch6:eq:values} gives
\[
V_{s,\mathrm{root}}=\sum_{i=1}^{K}q_{s,i}V_{s,\{i\}}
=\sum_{i=1}^{K}q_{s,i}A_{s,i}.
\]
\end{proof}

\begin{Thm:Theorem}[General prediction-error bound]\label{ch6:thm:master}
For any causal differentiable convex predictors and fixed positive step sizes,
Algorithm~\ref{ch6:alg:controller} is deadline feasible and, for every $\mathbf m\in\Delta_K$,
\begin{equation}\label{ch6:eq:master}
 \boxed{J_T^{\mathrm{on}}-J_T(\mathbf m)\le
 \sum_{I\in\mathcal{I}_K}M_I\left[
 \frac{d_I{D_\psi}(p_I^*,1/2)}{\eta_I}
 +\frac{\eta_I}{2\sigma_\psi}\sum_{s=1}^Ne_{s,I}^2\right].}
\end{equation}
With strictly positive error bounds $B_I$ and steps \eqref{ch6:eq:steps}, it also satisfies
\begin{equation}\label{ch6:eq:path}
 \Reg_T(\Delta_K)\le{C_\psi}\sqrt N
 \max_{1\le i\le K}\sum_{\substack{I\in\mathcal{I}_K\\i\in I}}B_I\sqrt{d_I}.
\end{equation}
\end{Thm:Theorem}
\begin{proof}
Proposition~\ref{ch6:prop:state} gives deadline feasibility.  For any
$\mathbf m\in\Delta_K$, Lemmas~\ref{ch6:lem:joint} and
\ref{ch6:lem:tree} give
\begin{align*}
J_T^{\mathrm{on}}-J_T(\mathbf m)
&\le\sum_{s=1}^{N}\sum_{I\in\mathcal I_K}
 M_I(p_{s,I}-p_I^*)g_{s,I}\\
&=\sum_{I\in\mathcal I_K}M_I
 \sum_{s=1}^{N}g_{s,I}(p_{s,I}-p_I^*).
\end{align*}
Since $M_I\ge0$, Lemma~\ref{ch6:lem:local} gives \eqref{ch6:eq:master}.
With \eqref{ch6:eq:steps} and $|e_{s,I}|\le B_I$,
\begin{align*}
J_T^{\mathrm{on}}-J_T(\mathbf m)
&\le C_\psi\sqrt N\sum_{I\in\mathcal I_K}
 M_IB_I\sqrt{d_I}\\
&=C_\psi\sqrt N\sum_{i=1}^{K}m_i
 \sum_{\substack{I\in\mathcal I_K\\i\in I}}B_I\sqrt{d_I}\\
&\le C_\psi\sqrt N
 \max_{1\le i\le K}\sum_{\substack{I\in\mathcal I_K\\i\in I}}
 B_I\sqrt{d_I},
\end{align*}
because $\sum_i m_i=1$.  Apply the last display to a minimizer of
$J_T$ over $\Delta_K$.
\end{proof}

\subsection{General upper bounds}
Consider two service ages $i<j$ in an interval $I$ containing $d_I$ ages.
Both marginal-cost coefficients include the storage terms for ages
$1,\ldots,i$. These common terms cancel when we subtract $A_{s,i}$ from
$A_{s,j}$, giving
\begin{equation}\label{ch6:eq:leaf-difference}
 A_{s,j}-A_{s,i}=w_s\left[
 \zeta_{s+j}^u-\zeta_{s+i}^u+\sum_{k=i+1}^{j}\zeta_{s+k}^x\right].
\end{equation}
The remaining sum contains $j-i$ storage terms. Since $j-i\le d_I-1$,
their total absolute contribution is at most $w_s(d_I-1)L_x$. The two
service terms contribute at most $2w_sL_u$. Each child-subtree value is an
average of the coefficients for its service ages, so the difference between
the two child values satisfies the same bound. Therefore, for the basic
controller ($F=0$),
\begin{equation}\label{ch6:eq:general-node-bound}
 |e_{s,I}|=|g_{s,I}|\le w_sG_I\le G_I,
 \qquad G_I=2L_u+(d_I-1)L_x.
\end{equation}

Define
\begin{equation}\label{ch6:eq:S}
 \Lambda_1=0,\qquad
 \Lambda_K=\sum_{j=0}^{\lceil\log_2K\rceil-1}
 \sqrt{\lceil K/2^j\rceil}\quad(K\ge2).
\end{equation}
At depth $j$ every interval has length at most $\lceil K/2^j\rceil$.
Consequently
\begin{equation}\label{ch6:eq:S-bound}
 \Lambda_K\le\sqrt{2K}\sum_{j\ge0}2^{-j/2}
 =2(1+\sqrt2)\sqrt K.
\end{equation}

\begin{Thm:Theorem}[Adversarial jointly convex costs]\label{ch6:thm:general}
Under \eqref{ch6:eq:support}--\eqref{ch6:eq:bounded-grad}, use $K=H$, $F=0$,
and \eqref{ch6:eq:steps} with $B_I=G_I$. The controller satisfies
\begin{align}
 \Reg_T(\Delta_H)
 &\le{C_\psi}\sqrt{T-1}
 \max_{1\le i\le H}\sum_{I\ni i}
 [2L_u+(d_I-1)L_x]\sqrt{d_I}\label{ch6:eq:general-path}\\
 &\le{C_\psi}[2L_u+(H-1)L_x]\sqrt{T-1}\,\Lambda_{H}.
 \label{ch6:eq:general-common}
\end{align}
In particular, for quadratic updates,
\begin{equation}\label{ch6:eq:main-upper}
 \boxed{\Reg_T(\Delta_H)\le
 \frac{2L_u+(H-1)L_x}{2}\sqrt{T-1}\,\Lambda_{H}
 =O\bigl(L_u\sqrt{HT}+L_xH^{3/2}\sqrt T\bigr).}
\end{equation}
\end{Thm:Theorem}
\begin{proof}
With $K=H$ and $N=T-1$, \eqref{ch6:eq:general-node-bound} gives
\[
B_I=G_I=2L_u+(d_I-1)L_x
\le 2L_u+(H-1)L_x=:G_*.
\]
For every leaf $i$, there is at most one interval containing $i$ at each
depth $j$, and its length is at most $\lceil H/2^j\rceil$.  Therefore
\begin{align*}
\sum_{\substack{I\in\mathcal I_{H}\\i\in I}}
G_I\sqrt{d_I}
&\le G_*\sum_{j=0}^{\lceil\log_2H\rceil-1}
\sqrt{\lceil H/2^j\rceil}\\
&=G_*\Lambda_{H}.
\end{align*}
Theorem~\ref{ch6:thm:master} gives
\eqref{ch6:eq:general-path}--\eqref{ch6:eq:general-common}.
For quadratic updates, $C_\psi=1/2$; \eqref{ch6:eq:S-bound} then gives
\eqref{ch6:eq:main-upper}.
\end{proof}

\begin{Thm:Corollary}[Consumption-only costs]\label{ch6:cor:action}
If $L_x=0$, the quadratic controller satisfies
\begin{equation}\label{ch6:eq:action-upper}
 \Reg_T(\Delta_H)\le L_u\sqrt{T-1}\,\Lambda_{H}
 \le2(1+\sqrt2)L_u\sqrt{H(T-1)}.
\end{equation}
\end{Thm:Corollary}
\begin{proof}
Set $L_x=0$ in \eqref{ch6:eq:general-common} and use
\eqref{ch6:eq:S-bound}. State independence follows from the Lipschitz bound
\eqref{ch6:eq:lipschitz} with $L_x=0$.
\end{proof}

For fixed $H,L_x,L_u$, the regret is $O(\sqrt T)$. With fixed positive
bounds and a growing deadline, $H^3=o(T)$ is sufficient for the general
upper bound to be sublinear, whereas $H=o(T)$ suffices when $L_x=0$.

\section{Deadline Selection under Nonnegative State Marginals}
\label{ch6:sec:positive}
This section adds a structural assumption on the joint costs $c_t$.
Assume that at every time
and every feasible pair the oracle can select a joint subgradient with
\begin{equation}\label{ch6:eq:positive}
 0\le\lambda\le\zeta^x\le\lambda+V,\qquad |\zeta^u|\le L_u,
 \qquad \lambda,V\ge0.
\end{equation}
The constant $\lambda$ is a lower bound on the marginal cost of holding
additional resource. The quantity $V$ specifies how far above $\lambda$
this marginal cost may rise. These same bounds apply at every round and
every feasible stored and consumed amount. The bounds are known,
although the cost functions themselves need not be known in advance.

\subsection{A comparison inequality for coupled costs}
\begin{Thm:Lemma}[Moving consumption earlier]\label{ch6:lem:earlier}
Let two feasible pairs satisfy $\bar x\le x$. Under \eqref{ch6:eq:positive},
\begin{equation}\label{ch6:eq:earlier}
 c_t(\bar x,\bar u)-c_t(x,u)
 \le-\lambda(x-\bar x)+L_u|\bar u-u|.
\end{equation}
\end{Thm:Lemma}
\begin{proof}
Use a bounded joint subgradient $(\bar\zeta^x,\bar\zeta^u)$ at
$(\bar x,\bar u)$ in \eqref{ch6:eq:support}, tested against $(x,u)$:
\begin{align*}
c_t(\bar x,\bar u)-c_t(x,u)
&\le\bar\zeta^x(\bar x-x)+\bar\zeta^u(\bar u-u)\\
&\le\lambda(\bar x-x)+|\bar\zeta^u|\,|\bar u-u|\\
&\le-\lambda(x-\bar x)+L_u|\bar u-u|,
\end{align*}
where the second line uses $\bar x-x\le0$ and
$\bar\zeta^x\ge\lambda$.
\end{proof}

To justify a shorter deadline, compare service at an age $i>K$ with
service at age one. If the age-$i$ service occurs by time $T$, moving one
unit changes consumption in two rounds. Lemma~\ref{ch6:lem:earlier} bounds
the contribution of these action changes to the cost increase by $2L_u$.
The move also removes $i-1\ge K$ rounds of storage, whose contribution to
the same bound is at most $-\lambda K$. This motivates choosing $K$ so
that $\lambda K\ge2L_u$ whenever the deadline is shortened. The next lemma
accounts separately for arrivals whose service extends beyond the cost
horizon. Choose
\begin{equation}\label{ch6:eq:K}
 K=\begin{cases}
 H,&\lambda=0,\\[1mm]
 \displaystyle\min\left\{H,
 \max\left\{1,\left\lceil\frac{2L_u}{\lambda}\right\rceil\right\}\right\},
 &\lambda>0.
 \end{cases}
\end{equation}
This choice also gives $\lambda(K-1)\le2L_u$, so the part of the local
feedback bound \eqref{ch6:eq:positive-B} proportional to $\lambda$ is at
most $2L_u$. When $K<H$, it satisfies
$\lambda K\ge2L_u$ as required for the preceding cost comparison.
For $K=1$ there are no split decisions to learn.

\begin{Thm:Lemma}[Cost of shortening the comparator's deadline]\label{ch6:lem:reduce}
For any $\mathbf m\in\Delta_{H}$, move all its fractions at ages $i>K$ to age one,
producing $\bar{\mathbf m}\in\Delta_K$. Then
\begin{equation}\label{ch6:eq:reduce}
 J_T(\bar{\mathbf m})-J_T(\mathbf m)\le L_uK\,\ind{K<H}.
\end{equation}
No modification is needed when $K=H$.
\end{Thm:Lemma}
\begin{proof}
For $i>K$ and $s\le T-1$, put
\[
v_{s,i}=w_sm_i,\qquad h_s=T-s.
\]
Replacing age-$i$ service of this amount by age-one service changes the
actions within the cost horizon only at $s+1$ and, when $s+i\le T$, at $s+i$.
Its action and state contributions are therefore
\begin{align*}
\sum_{t=1}^{T}|\Delta u_t^{(s,i)}|
&=v_{s,i}\bigl(1+\ind{i\le h_s}\bigr),\\
\sum_{t=1}^{T}\bigl(x_t(\mathbf m)-x_t(\bar{\mathbf m})\bigr)^{(s,i)}
&=v_{s,i}\sum_{r=2}^{\min\{i,h_s\}}1\\
&=v_{s,i}\bigl(\min\{i,h_s\}-1\bigr).
\end{align*}
In particular, $x_t(\bar{\mathbf m})\le x_t(\mathbf m)$ for every $t$.
State reductions add linearly, while the action changes satisfy
\[
\sum_{t=1}^{T}|u_t(\bar{\mathbf m})-u_t(\mathbf m)|
\le\sum_{s=1}^{T-1}\sum_{i=K+1}^{H}
v_{s,i}\bigl(1+\ind{i\le h_s}\bigr).
\]
After summing Lemma~\ref{ch6:lem:earlier} over $t$,
\begin{align}
J_T(\bar{\mathbf m})-J_T(\mathbf m)
&\le\sum_{s=1}^{T-1}\sum_{i=K+1}^{H}v_{s,i}
\left[
L_u\bigl(1+\ind{i\le h_s}\bigr)
-\lambda\bigl(\min\{i,h_s\}-1\bigr)
\right].
\label{ch6:eq:moved-net}
\end{align}
If $K=H$, the double sum is empty.  Otherwise,
\eqref{ch6:eq:K} gives $\lambda K\ge2L_u$, so for $i>K$,
\[
L_u(1+\ind{i\le h_s})-\lambda(\min\{i,h_s\}-1)
\le
\begin{cases}
2L_u-\lambda K\le0, & h_s\ge K+1,\\
L_u-\lambda(h_s-1)\le L_u, & 1\le h_s\le K.
\end{cases}
\]
Consequently,
\begin{align*}
J_T(\bar{\mathbf m})-J_T(\mathbf m)
&\le L_u\sum_{\substack{1\le s\le T-1\\1\le T-s\le K}}
\sum_{i=K+1}^{H}w_sm_i\\
&\le L_u\sum_{s=T-K}^{T-1}1
\le L_uK.
\end{align*}
\end{proof}

For arrivals near the end of the cost horizon, moving service earlier can
add a positive action term to the cost bound at or before $T$, while part
of the compensating storage reduction occurs after $T$ and is excluded
from the objective. When $K<H$, only arrivals $s=T-K,\ldots,T-1$ can
contribute a positive excess in the preceding bound. Their total size is at
most $K$, and each unit contributes at most $L_u$, giving the term $L_uK$ in
\eqref{ch6:eq:reduce}. This comparison allows the controller to use deadline
$K$ while its regret is still measured against fixed policies with deadline
$H$. The comparison bound can be replaced by zero when $K=H$, or when
every nonzero arrival's full original lifetime lies inside the cost horizon.

\subsection{Regret bound for the chosen deadline}
For the basic controller, \eqref{ch6:eq:leaf-difference} and
\eqref{ch6:eq:positive} imply
\begin{equation}\label{ch6:eq:positive-node}
 |g_{s,I}|\le w_s[2L_u+(\lambda+V)(d_I-1)].
\end{equation}
A common bound is
\begin{equation}\label{ch6:eq:positive-B}
 B=2L_u+(\lambda+V)(K-1)\le4L_u+(K-1)V.
\end{equation}

\begin{Thm:Theorem}[Regret bound with the chosen deadline]\label{ch6:thm:positive}
Assume \eqref{ch6:eq:positive}, use $K$ in \eqref{ch6:eq:K}, $F=0$, and steps
\eqref{ch6:eq:steps} with the common bound \eqref{ch6:eq:positive-B}. Then
\begin{equation}\label{ch6:eq:positive-upper}
 \boxed{\Reg_T(\Delta_H)\le
 {C_\psi}[2L_u+(\lambda+V)(K-1)]\sqrt{T-1}\,\Lambda_K
 +L_uK\,\ind{K<H}.}
\end{equation}
In particular,
\begin{equation}\label{ch6:eq:positive-rate}
 \Reg_T(\Delta_H)=O\bigl((L_u+KV)\sqrt{KT}+L_uK\bigr).
\end{equation}
\end{Thm:Theorem}
\begin{proof}
Let $\mathbf m^*$ minimize $J_T$ over $\Delta_{H}$, and let
$\bar{\mathbf m}^*\in\Delta_K$ be its reduction from
Lemma~\ref{ch6:lem:reduce}.  With $B$ from \eqref{ch6:eq:positive-B},
\eqref{ch6:eq:positive-node} gives $|g_{s,I}|\le B$ at every node.
Theorem~\ref{ch6:thm:master} and the same depth calculation as in
Theorem~\ref{ch6:thm:general} give
\[
J_T^{\mathrm{on}}-J_T(\bar{\mathbf m}^*)
\le C_\psi B\sqrt{T-1}\,\Lambda_K.
\]
Therefore
\begin{align*}
\Reg_T(\Delta_H)
&=J_T^{\mathrm{on}}-J_T(\mathbf m^*)\\
&=J_T^{\mathrm{on}}-J_T(\bar{\mathbf m}^*)
  +J_T(\bar{\mathbf m}^*)-J_T(\mathbf m^*)\\
&\le C_\psi[2L_u+(\lambda+V)(K-1)]
  \sqrt{T-1}\,\Lambda_K
  +L_uK\ind{K<H}.
\end{align*}
Equations \eqref{ch6:eq:S-bound} and \eqref{ch6:eq:positive-B} give
\eqref{ch6:eq:positive-rate}.
Using the node-dependent bounds in \eqref{ch6:eq:positive-node} directly
in \eqref{ch6:eq:path} gives the sharper path expression.
\end{proof}

\begin{Thm:Corollary}[Constant nonnegative state marginal]\label{ch6:cor:constant}
If the oracle selects $\zeta^x=\lambda\ge0$ at every time and feasible pair,
then the quadratic controller satisfies
\begin{equation}\label{ch6:eq:constant}
 \Reg_T(\Delta_H)\le2L_u\sqrt{T-1}\,\Lambda_K
 +L_uK\,\ind{K<H}
 =O(L_u\sqrt{KT}+L_uK).
\end{equation}
The implied constant is independent of $\lambda$, and this is at most
$O(L_u\sqrt{HT})$.
\end{Thm:Corollary}
\begin{proof}
With $V=0$, $C_\psi=1/2$, and \eqref{ch6:eq:K},
\[
C_\psi[2L_u+\lambda(K-1)]
\le\frac12(2L_u+2L_u)=2L_u.
\]
Apply Theorem~\ref{ch6:thm:positive}, \eqref{ch6:eq:S-bound}, and
$K\le\sqrt{KT}\le\sqrt{HT}$.
\end{proof}

When $\lambda>0$ and $L_u$ are fixed, the deadline $K$ is bounded
independently of $H$. For example, $\lambda=L_u=1$ gives
$K=\min\{H,2\}$. When only $\zeta^x\ge0$ is known and no positive lower
bound is available, \eqref{ch6:eq:K} gives $K=H$: long zero-cost holding
periods cannot be excluded by this argument.
Section~\ref{ch6:subsec:spikes} shows why nonnegativity alone does not give a
bound independent of the magnitude of adversarial state marginals.

\section{Causal Prediction and Known State Costs}
\label{ch6:sec:prediction}
Forecasts of future costs can help the controller decide how much resource
to consume in the earlier interval and how much to leave for the later
interval. A forecast may use previously revealed costs, side information,
or a known part of the cost function. It must be available when the split
is chosen. Setting the prediction term $F$ to zero gives the basic
controller.

\subsection{A hypothetical trajectory used for prediction}
To use a forecast when choosing a split, we first give it a simple
hypothetical consumption schedule. Part of the arrival is consumed at the
start of the earlier interval, and the rest at the start of the later
interval. Varying that fraction gives a predicted cost for each candidate
split. The actual consumption times are still determined by later decisions
in the tree.

Fix $I=[a,b]$, let $r=r_I$ be its right-child starting age, and consider
arrival $s$. At the end of $s+a-1$, choose a known background storage level
$B_{s,I}^{\mathrm{bg}}\in[0,K-w_s]$; zero is always allowed.
A hypothetical allocation of size $w_s$ serves fraction $p$ at age
$a$ and fraction $1-p$ at age $r$. In the hypothetical states, actions,
and forecast gradients below, we suppress the fixed indices $(s,I)$;
age $k$ corresponds to physical round $s+k$. For $k=a,\ldots,r$, define
\begin{align}
 \widehat x_k(p)&=B_{s,I}^{\mathrm{bg}}
 +w_s\bigl[p\ind{k=a}+(1-p)\bigr],\label{ch6:eq:virtual-x}\\
 \widehat u_k(p)&=w_s\bigl[p\ind{k=a}
 +(1-p)\ind{k=r}\bigr].\label{ch6:eq:virtual-u}
\end{align}
Every such pair is affine in $p$ and belongs to $\mathcal{D}_K$.
The background is retained and receives no service in the hypothetical schedule during this
interval. The size is the full $w_s$, even if the node receives less actual
mass, consistently with the normalization of \eqref{ch6:eq:g}.

For each $k=a,\ldots,r$ with $s+k\le T$, let
$\widehat c_{s,I,k}(x,u)$ be a causal differentiable jointly convex forecast
for $c_{s+k}$. Put
\begin{equation}\label{ch6:eq:joint-predictor}
 F_{s,I}(p)=
 \sum_{k=a}^{\min\{r,T-s\}}
 \widehat c_{s,I,k}(\widehat x_k(p),\widehat u_k(p)),
\end{equation}
with an empty sum equal to zero. Then \eqref{ch6:eq:play} retains these full
nonlinear functions inside a one-dimensional convex optimization problem.

To identify the predicted gradient, evaluate the forecast gradient
$(\widehat\zeta_k^x,\widehat\zeta_k^u)$ at the hypothetical state--action pair for the chosen
$p_{s,I}$. The chain rule gives
\begin{equation}\label{ch6:eq:predicted-g}
 \widehat g_{s,I}=w_s\left[
 \ind{s+a\le T}\widehat\zeta_a^u
 -\ind{s+r\le T}\widehat\zeta_r^u
 -\sum_{k=a+1}^{\min\{r,T-s\}}\widehat\zeta_k^x
 \right].
\end{equation}
Omitted forecast terms are zero. Equation~\eqref{ch6:eq:correct}, centered at
the old reference, performs the eventual correction using the actual
$g_{s,I}$.

The forecast evaluates costs using a simplified consumption schedule and
hypothetical storage levels. The actual controller may consume resource
at different ages and reach different storage levels, so the resulting
marginal comparisons can differ even when the forecast cost functions
are exact. Thus $e_{s,I}=g_{s,I}-\widehat g_{s,I}$ includes both errors in
the cost forecasts and differences caused by the hypothetical schedule.
The regret bound in \eqref{ch6:eq:master} depends on the sum of these
squared feedback errors.

\subsection{Error bounds and worst-case guarantees}
\begin{Thm:Proposition}[Forecasts with bounded gradients]\label{ch6:prop:forecast}
If every forecast gradient obeys
$|\widehat\zeta^x|\le L_x$ and $|\widehat\zeta^u|\le L_u$ on its hypothetical
pairs, then
\begin{equation}\label{ch6:eq:forecast-general}
 |e_{s,I}|\le2w_s[2L_u+(d_I-1)L_x].
\end{equation}
Thus steps based on twice the basic feedback bound preserve the general
regret order of Theorem~\ref{ch6:thm:general}.
If instead actual subgradients and forecast gradients both satisfy the bounds in
\eqref{ch6:eq:positive}, then
\begin{equation}\label{ch6:eq:forecast-positive}
 |e_{s,I}|\le w_s[4L_u+(\lambda+V)(d_I-1)].
\end{equation}
Together with \eqref{ch6:eq:K}, this preserves the rate
$O((L_u+KV)\sqrt{KT}+L_uK)$ in Theorem~\ref{ch6:thm:positive}.
\end{Thm:Proposition}
\begin{proof}
Write $g_{s,I}=g_{s,I}^u+g_{s,I}^x$ and
$\widehat g_{s,I}=\widehat g_{s,I}^u+\widehat g_{s,I}^x$ for the
action- and state-marginal terms.  Equation~\eqref{ch6:eq:predicted-g}
contains at most two action terms and at most $d_I-1$ state terms, so
\[
|\widehat g_{s,I}|
\le w_s\bigl[2L_u+(d_I-1)L_x\bigr].
\]
Together with \eqref{ch6:eq:general-node-bound},
\begin{align*}
|e_{s,I}|
&=|g_{s,I}-\widehat g_{s,I}|\\
&\le|g_{s,I}|+|\widehat g_{s,I}|\\
&\le2w_s\bigl[2L_u+(d_I-1)L_x\bigr].
\end{align*}

Under \eqref{ch6:eq:positive}, the subtree values are weighted
averages of leaf values. For fixed $(s,I)$, let $\pi_i^{\mathrm{L}}$ and
$\pi_i^{\mathrm{R}}$ be the controller's final conditional leaf fractions
in the left and right children, respectively, with their dependence on
$(s,I)$ suppressed. These probability weights give
\[
g_{s,I}
=\sum_{i\in I_{\mathrm{L}}}\sum_{j\in I_{\mathrm{R}}}\pi_i^{\mathrm{L}}\pi_j^{\mathrm{R}}(A_{s,i}-A_{s,j}).
\]
For every $i\in I_{\mathrm{L}}$ and $j\in I_{\mathrm{R}}$,
\begin{align*}
A_{s,i}-A_{s,j}
&=w_s\left[
\zeta_{s+i}^u-\zeta_{s+j}^u
-\sum_{k=i+1}^{j}\zeta_{s+k}^x\right],\\
\left|w_s(\zeta_{s+i}^u-\zeta_{s+j}^u)\right|
&\le2L_uw_s,\\
-w_s(\lambda+V)(d_I-1)
&\le-w_s\sum_{k=i+1}^{j}\zeta_{s+k}^x\le0.
\end{align*}
Averaging and applying \eqref{ch6:eq:predicted-g} to the forecast gives
\[
|g_{s,I}^u|,|\widehat g_{s,I}^u|\le2L_uw_s,
\qquad
g_{s,I}^x,\widehat g_{s,I}^x
\in[-w_s(\lambda+V)(d_I-1),0].
\]
Hence
\begin{align*}
|e_{s,I}|
&\le|g_{s,I}^u-\widehat g_{s,I}^u|
  +|g_{s,I}^x-\widehat g_{s,I}^x|\\
&\le w_s\bigl[4L_u+(\lambda+V)(d_I-1)\bigr].
\end{align*}
Finally,
\[
4L_u+(\lambda+V)(K-1)
\le6L_u+(K-1)V
\]
by $\lambda(K-1)\le2L_u$.  Theorem~\ref{ch6:thm:master},
\eqref{ch6:eq:S-bound}, and Lemma~\ref{ch6:lem:reduce} therefore give
$O((L_u+KV)\sqrt{KT}+L_uK)$.
\end{proof}

At each node, the error contribution in \eqref{ch6:eq:master} is proportional
to $\sum_s(g_{s,I}-\widehat g_{s,I})^2$. Accurate predictions of the marginal
feedback can make this contribution small. Uniform gradient bounds alone
do not ensure that improvement: they allow the actual and predicted
marginals to differ by twice the basic feedback bound, which explains the
larger general worst-case constant in Proposition~\ref{ch6:prop:forecast}.

These forecast bounds assume differentiable forecast functions and exact
solutions of the split-selection and reference-update minimizations. The observed costs themselves may be
nonsmooth.

\subsection{A fixed known state-cost component}
\begin{Thm:Corollary}[Known state cost]\label{ch6:cor:known-state}
Suppose the joint cost has the structured form
\begin{equation}\label{ch6:eq:known-state-cost}
 c_t(x,u)=\ell_t(u)+f(x),
\end{equation}
where each $\ell_t$ is adversarial and convex with available subgradients
bounded by $L_u>0$ on $[0,H]$, while $f$ is convex, differentiable, fixed,
and known before the first decision. Assume
$0\le\lambda\le f'(x)\le\lambda+V$ on $[0,H]$ and choose $K$ by
\eqref{ch6:eq:K}. Then the basic controller satisfies
Theorem~\ref{ch6:thm:positive} with the same constants.

When the storage cost $f$ is known, we can evaluate it directly when
choosing the split. In the hypothetical schedule, consuming fraction $p$
at age $a$ leaves $w_s(1-p)$ from this arrival in storage until age $r_I$.
Use the causal joint forecasts $\widehat c_{s,I,k}(x,u)=f(x)$ in
\eqref{ch6:eq:joint-predictor}, the hypothetical states \eqref{ch6:eq:virtual-x}, and set
\[
 \gamma_{s,I}=\max\{0,\min\{r_I,T-s\}-a\}.
\]
The quantity $\gamma_{s,I}$ counts the subsequent rounds through age
$r_I$ that fall within the cost horizon. It therefore determines how
many times the storage cost of this remaining amount enters the
prediction step.
After dropping a term independent of $p$, the prediction step becomes
\begin{equation}\label{ch6:eq:known-state-prox}
 \boxed{p_{s,I}=\argmin_{p\in[0,1]}
 \left\{{D_\psi}(p,z_{s,I})+
 \eta_I\gamma_{s,I}
 f\bigl(B_{s,I}^{\mathrm{bg}}+w_s(1-p)\bigr)\right\}.}
\end{equation}
Keep the reference correction \eqref{ch6:eq:correct} unchanged and use the
step sizes \eqref{ch6:eq:steps} with the bound \eqref{ch6:eq:positive-B}.
For either regularizer, this controller satisfies
\eqref{ch6:eq:positive-upper}. The prediction-error bound
\eqref{ch6:eq:master} also applies.
In particular, $f(x)=\lambda x+b$ with $\lambda\ge0$ recovers
Corollary~\ref{ch6:cor:constant} for the quadratic controller, including its
constant independent of $\lambda$.
\end{Thm:Corollary}
\begin{proof}
The joint subgradient $(f'(x_t),\zeta_t^u)$, with
$\zeta_t^u\in\partial\ell_t(u_t)$ and $|\zeta_t^u|\le L_u$, satisfies
\eqref{ch6:eq:positive}.  The basic-controller claim follows from
Theorem~\ref{ch6:thm:positive}.

For the forecast $\widehat c_{s,I,k}(x,u)=f(x)$,
\eqref{ch6:eq:virtual-x} gives
\[
\widehat x_a(p)=B_{s,I}^{\mathrm{bg}}+w_s,
\qquad
\widehat x_k(p)=B_{s,I}^{\mathrm{bg}}+w_s(1-p)
\quad(a<k\le r_I).
\]
Consequently,
\begin{align*}
F_{s,I}(p)
&=\ind{s+a\le T}f(B_{s,I}^{\mathrm{bg}}+w_s)\\
&\quad+\gamma_{s,I}
 f(B_{s,I}^{\mathrm{bg}}+w_s(1-p)).
\end{align*}
The first term is independent of $p$, and substitution in
\eqref{ch6:eq:play} gives \eqref{ch6:eq:known-state-prox}.

Here $\widehat g_{s,I}^u=0$, so the component bounds in the proof of
Proposition~\ref{ch6:prop:forecast} give
\begin{align*}
|e_{s,I}|
&\le|g_{s,I}^u|
  +|g_{s,I}^x-\widehat g_{s,I}^x|\\
&\le w_s[2L_u+(\lambda+V)(d_I-1)].
\end{align*}
Theorem~\ref{ch6:thm:master} and Lemma~\ref{ch6:lem:reduce} then give
\eqref{ch6:eq:positive-upper}.
For $f(x)=\lambda x+b$, take $V=0$ and apply the quadratic specialization
in Corollary~\ref{ch6:cor:constant}.
\end{proof}

If a known convex
$f$ is not nondecreasing, \eqref{ch6:eq:known-state-cost} still falls under
Theorem~\ref{ch6:thm:general} whenever $|f'|\le L_x$, and the same
predictor is covered by Proposition~\ref{ch6:prop:forecast} with $K=H$.
The refinement that can shorten the deadline then need not apply.

\section{Implementation and Deadline Dependence}
\label{ch6:sec:implementation}
At a node of length $d_I$, store its $d_I$ reference values and the outstanding
splits and child values needed to complete feedback. Each node starts at
most one new decision for an arrival per round, and its feedback is complete
within $d_I$ rounds. It therefore needs at most $d_I$ unfinished feedback
records. Each record stores a constant number of split, routing, and
child-value entries, so the node uses $O(d_I)$ memory in total. Resource
reserved for a later interval remains stored until that interval begins.
Since intervals at a fixed depth are disjoint and have total length at most $K$,
\begin{equation}\label{ch6:eq:memory}
 \sum_{I\in\mathcal{I}_K}d_I=O(K\log(2K)).
\end{equation}
Summing over the tree gives $O(K\log(2K))$ memory. Completed records can
be discarded, so this requirement does not grow with $T$.

Compute the leaf values using a scalar prefix sum
$P_t=\sum_{\tau=1}^t\zeta_\tau^x$ and its stored values at active arrival times:
\[
 A_{s,i}=w_s(\zeta_{s+i}^u+P_{s+i}-P_s).
\]
Each node has at most one new split decision for an arrival and one completed feedback value per
physical round. A direct implementation uses $O(K\log(2K))$ arithmetic per
round. This bound excludes evaluating or solving arbitrary forecast problems;
\eqref{ch6:eq:joint-predictor} may require several predicted cost evaluations.

Both regularizers use the same state, routing, feedback timing, and tree
identity. Quadratic updates are additive in a fraction; entropic updates are
additive in its log odds. They generally produce different trajectories.
Nonlinear prediction is available with either regularizer.

The deadline affects this controller's regret bound in two ways. First,
the root must wait $K$ rounds for complete feedback and therefore uses
$K$ separate learner copies. Combining their regret bounds introduces a
factor of $\sqrt K$. Second, changing the consumption age can change as
many as $K-1$ storage-cost terms, so the feedback magnitude grows as
$O(L_u+KL_x)$. The shorter intervals farther down the tree add only a
constant factor. Together, these effects give the bound
\[
 O\bigl((L_u+KL_x)\sqrt{KT}\bigr)
 =O\bigl(L_u\sqrt{KT}+L_xK^{3/2}\sqrt T\bigr).
\]
With $K=H$, these are the service and storage contributions in
\eqref{ch6:eq:main-upper}.
The storage-cost lower bound in the next section matches the state contribution.

\section{Lower Bounds and Minimax Implications}
\label{ch6:sec:lower}
General optimization with memory already exhibits an
$\Omega(Lh^{3/2}\sqrt T)$ lower bound for a suitable memory-$h$
loss class \cite[Appendix D]{kumar2023unbounded}. The construction here
establishes the storage-dependent rate within the admissible scalar
system, against its fixed age policies and for all causal deadline-feasible
controllers; it is not an automatic consequence of that broader class.

All lower bounds below use oblivious instances: arrivals and cost functions
are fixed before the first decision and are independent of the controller's internal
randomization. Random signs are used only to construct a distribution over
such fixed instances. The expectation in the final statements for fixed
input sequences is over the controller alone.

Appendix~\ref{app:signs} supplies the independent-sign moment bound
and the averaging argument used below. In particular,
$\E[|S_n|]\ge\sqrt{n/3}$ gives the constants used here.

\subsection{The action-marginal lower bound}
\begin{Thm:Theorem}[Consumption-only lower bound]\label{ch6:thm:lower-action}
For $T\ge2$, $H\ge2$, and $L_u>0$, every causal deadline-feasible controller
has a deterministic oblivious instance with consumption-only jointly affine
costs such that
\begin{equation}\label{ch6:eq:lower-action}
 \E[\Reg_T(\Delta_H)]\ge\frac{L_u}{2\sqrt3}\sqrt{T-1}.
\end{equation}
\end{Thm:Theorem}
\begin{proof}
Set $w_s=1$ for odd $s$ and $w_s=0$ for even $s$. Draw independent signs
$\xi_2,\ldots,\xi_T$ and set
$c_1(x,u)=0$, $c_t(x,u)=L_u\xi_tu$ for $t\ge2$.
Let $\mathcal F_{t-1}^{\mathcal A}$ contain the controller's random seed
and all observations available before its round-$t$ action.  Then
$u_t$ is $\mathcal F_{t-1}^{\mathcal A}$-measurable and the fresh
obliviously sampled sign satisfies
$\E[\xi_t\mid\mathcal F_{t-1}^{\mathcal A}]=0$.  Hence
\[
\E[\xi_tu_t\mid\mathcal F_{t-1}^{\mathcal A}]
=u_t\E[\xi_t\mid\mathcal F_{t-1}^{\mathcal A}]=0,
\qquad
\E[J_T^{\mathrm{on}}]=L_u\sum_{t=2}^{T}\E[\xi_tu_t]=0.
\]
For the policies $\mathbf e_1,\mathbf e_2$, put
\[
C_1=L_u\sum_{t=2}^{T}\xi_tw_{t-1},
\qquad
C_2=L_u\sum_{t=2}^{T}\xi_tw_{t-2}.
\]
Then $\E[C_1]=\E[C_2]=0$ and
\[
w_{t-1}-w_{t-2}=(-1)^t\in\{-1,1\},
\qquad
C_1-C_2=L_u\sum_{t=2}^{T}(-1)^t\xi_t
\stackrel{\mathrm{d}}{=}L_uS_{T-1}.
\]
Using $\min\{a,b\}=(a+b-|a-b|)/2$,
\[
\E[\min\{C_1,C_2\}]
=-\frac12\E[|C_1-C_2|]
=-\frac{L_u}{2}\E[|S_{T-1}|].
\]
Since $\min_{\mathbf m\in\Delta_H}J_T(\mathbf m)\le\min\{C_1,C_2\}$,
\begin{align*}
\E[\Reg_T(\Delta_H)]
&\ge \E[J_T^{\mathrm{on}}]-\E[\min\{C_1,C_2\}]\\
&=\frac{L_u}{2}\E[|S_{T-1}|]\\
&\ge\frac{L_u}{2\sqrt3}\sqrt{T-1},
\end{align*}
by Lemma~\ref{ch6:lem:signs}.  Averaging over the finite sign family fixes
a deterministic sign sequence with at least this expected regret for the
given randomized controller.
\end{proof}

\subsection{A matching state-marginal lower bound}
\begin{Thm:Theorem}[Lower bound for signed linear storage costs]\label{ch6:thm:lower-state}
For $H\ge3$, $T\ge4$, and $L_x>0$, every causal deadline-feasible controller
has a deterministic oblivious instance with costs $c_t(x,u)=a_tx$,
$|a_t|\le L_x$, such that
\begin{equation}\label{ch6:eq:lower-state}
 \boxed{\E[\Reg_T(\Delta_H)]\ge
 \frac{L_x}{256\sqrt3}\,H^{3/2}\sqrt T.}
\end{equation}
The state-marginal term in Theorem~\ref{ch6:thm:general} is therefore sharp up
to universal constants in this range.
\end{Thm:Theorem}
\begin{proof}
Let
\begin{equation}\label{ch6:eq:block-parameters}
 b=\min\{\lfloor H/3\rfloor,\lfloor T/4\rfloor\},\qquad
 n=\lfloor T/(4b)\rfloor.
\end{equation}
We have $b,n\ge1$ and $4bn\le T$.  Use $n$ consecutive blocks of $4b$ rounds, with zero
costs and arrivals in any remaining rounds. In block $j$, draw a fresh
independent symmetric sign $\varepsilon_j$. At local time $r$ within the
block, specify the arrivals and coefficients in Table~\ref{tab:deadline-block-phases};
empty rows are omitted.
\begin{table}[tbp]
\centering
\caption{Local rounds in the deadline-allocation lower-bound construction.}
\label{tab:deadline-block-phases}
\begin{tabular}{@{}lll@{}}
\toprule
Local rounds $r$ & Arrival $w_t$ & Coefficient $a_t$ in $c_t(x,u)=a_tx$\\
\midrule
$1,\ldots,b$ & $1$ & $0$\\
$b+1$ & $0$ & $0$\\
$b+2,\ldots,2b+1$ & $0$ & $L_x/2$\\
$2b+2,\ldots,3b+1$ & $0$ & $-L_x(1+\varepsilon_j)/2$\\
$3b+2,\ldots,4b$ & $0$ & $0$\\
\bottomrule
\end{tabular}
\end{table}
The coefficient in the random phase is either $0$ throughout or $-L_x$
throughout. All signs can be sampled before the interaction.

Let $Y_j$ be the controller's state at the beginning of the first random-phase
round. It is determined before $\varepsilon_j$ is first revealed.
There are no arrivals throughout the positive-cost phase or the random
phase, so the state is nonincreasing there. Every positive-cost state is
at least $Y_j$, and every random-phase state is at most $Y_j$.
This remains true if resource from earlier blocks is still present.
Therefore the block's online cost satisfies the pathwise bound
\begin{align}
 C_j^{\mathrm{on}}
 &\ge\frac{L_x}{2}bY_j
 -\frac{L_x}{2}(1+\varepsilon_j)bY_j\nonumber\\
 &=-\frac{L_x}{2}b\varepsilon_jY_j.
 \label{ch6:eq:block-online}
\end{align}
Let $\mathcal F_j$ be the history, including the controller's random seed,
immediately before block $j$'s random phase.  Then $Y_j$ is
$\mathcal F_j$-measurable and
$\E[\varepsilon_j\mid\mathcal F_j]=0$.  Thus
\[
\E[C_j^{\mathrm{on}}\mid\mathcal F_j]
\ge-\frac{L_x}{2}bY_j\E[\varepsilon_j\mid\mathcal F_j]=0,
\qquad
\E[J_T^{\mathrm{on}}]=\sum_{j=1}^{n}\E[C_j^{\mathrm{on}}]\ge0.
\]
In particular, learning the sign during that phase cannot invalidate the
pathwise lower bound.

Consider the fixed comparators $\mathbf e_1$ and $\mathbf e_{3b}$. Both belong to $\Delta_H$.
The immediate policy has consumed every new arrival by local round $b+1$,
and has zero block cost. The delayed policy holds all $b$ new units at the
beginning of every positive and random cost round: its first service occurs
at $3b+1$, after that round's state cost. Its last service occurs at $4b$,
so it carries nothing into the next block. Its block cost is exactly
\begin{align}
C_j(\mathbf e_{3b})
&=\frac{L_x}{2}\,b\cdot b
-\frac{L_x}{2}(1+\varepsilon_j)b\cdot b\nonumber\\
&=-\frac{L_x}{2}b^2\varepsilon_j.
\label{ch6:eq:block-comparator}
\end{align}
Consequently, with $S_n=\sum_j\varepsilon_j$,
\begin{align}
 \E[\Reg_T(\Delta_H)]
 &\ge \E[J_T^{\mathrm{on}}]
 -\E\left[\min\!\left\{0,-\frac{L_xb^2}{2}S_n\right\}\right]\nonumber\\
 &\ge-\E\left[\min\!\left\{0,-\frac{L_xb^2}{2}S_n\right\}\right]\nonumber\\
 &=\frac{L_xb^2}{4}\E[|S_n|]
 \ge\frac{L_xb^2}{4\sqrt3}\sqrt n.
 \label{ch6:eq:block-lower}
\end{align}
Since $H/3\ge1$ and $T/4\ge1$,
$\lfloor H/3\rfloor\ge H/8$ and
$\lfloor T/4\rfloor\ge T/8\ge H/8$, where the last inequality uses $H\le T$. Thus
\[
b\ge\frac{H}{8}.
\]
Also, $\lfloor q\rfloor\ge q/2$ for $q\ge1$ gives
\[
n=\left\lfloor\frac{T}{4b}\right\rfloor
\ge\frac{T}{8b}.
\]
Thus
\begin{align*}
\frac{L_xb^2}{4\sqrt3}\sqrt n
&\ge\frac{L_x}{8\sqrt6}b^{3/2}\sqrt T\\
&\ge\frac{L_x}{256\sqrt3}H^{3/2}\sqrt T.
\end{align*}
Averaging over the finite sign family fixes a deterministic oblivious
instance with the same expected-regret lower bound for the given controller.
\end{proof}

\begin{Thm:Remark}[Small deadlines and the sign of the marginal]
The restriction $H\ge3$ is only for the block construction.
For $H=2$, $T\ge4$, let $w_t=1$ and use independent coefficients
$a_t=L_x\xi_t$ for $t\ge3$, with earlier costs zero.
Every controller again has zero expected cost, while the states of
$\mathbf e_1,\mathbf e_2$ differ by one at each such time. Their comparison gives
$L_x\sqrt{T-2}/(2\sqrt3)$, which is the same order as
$L_xH^{3/2}\sqrt T$ for $H=2$.
For $T=2$, all states within the cost horizon are predetermined, so no positive
state-only regret lower bound is asserted. For $H=1$ regret is zero.
The matching state lower bound allows increasing storage to either raise
or lower the current cost. The construction uses negative storage slopes
during some rounds. Its costs can nevertheless be made nonnegative by
adding constants that do not depend on the controller's decisions.
These constants leave regret and the storage slopes unchanged.
\end{Thm:Remark}

\subsection{Why nonnegative state marginal costs still need a magnitude bound}
\label{ch6:subsec:spikes}
Nonnegative state marginals can still produce unpredictable storage
charges. The next construction shows that, when service also has a cost,
no regret bound that is uniformly sublinear in $T$ over this class can be
independent of the largest allowed state marginal.

\begin{Thm:Proposition}[Unpredictable nonnegative spikes]\label{ch6:prop:spikes}
For every integer $n\ge2$ and $L_u>0$, there is a distribution of oblivious
instances with $H=2$, $T=3n$, nonnegative jointly affine costs, and
$0\le\zeta^x\le nL_u$, $0\le\zeta^u\le L_u$, such that every causal
controller has expected regret at least
\begin{equation}\label{ch6:eq:spikes}
 nL_u(1-1/n)^n\ge nL_u/4.
\end{equation}
A deterministic instance with at least this expected regret therefore exists
for every randomized controller.
\end{Thm:Proposition}
\begin{proof}
Use $n$ independent three-round cycles. One unit arrives in the first round;
there are no other arrivals in the cycle. Set the first cost to zero, the
second to $c_t(x,u)=L_uu$, and the third to $c_t(x,u)=A_jx$, where
$A_j=nL_u$ with probability $1/n$ and $A_j=0$ otherwise.
These coefficients are sampled in advance and only revealed at their rounds.
The deadline forces all cycle resource to be consumed in its second and
third rounds, and all earlier resource is already gone.
If the controller serves fraction $p_j$ in the second round, its cycle cost
is $C_j^{\mathrm{on}}=L_up_j+A_j(1-p_j)$.
Since $p_j$ is measurable before $A_j$ is revealed and
$\E[A_j]=(1/n)nL_u=L_u$,
\begin{align*}
\E[C_j^{\mathrm{on}}\mid p_j]
&=L_up_j+\E[A_j\mid p_j](1-p_j)=L_u,\\
\E[J_T^{\mathrm{on}}]&=\sum_{j=1}^{n}\E[C_j^{\mathrm{on}}]=nL_u.
\end{align*}
For a fixed comparator $(p,1-p)^\top\in\Delta_2$,
\[
J_T(p)=pnL_u+(1-p)\sum_{j=1}^{n}A_j.
\]
Thus
\[
\min_{p\in[0,1]}J_T(p)
=\min\left\{nL_u,\sum_{j=1}^{n}A_j\right\}
=nL_u\ind{\max_jA_j>0}.
\]
Therefore
\begin{align*}
\E[\Reg_T(\Delta_2)]
&=nL_u\mathbb P(\max_jA_j=0)\\
&=nL_u\left(1-\frac1n\right)^n
\ge\frac{nL_u}{4},
\end{align*}
where $(1-1/n)^n\ge(1-1/2)^2=1/4$ for $n\ge2$.
Averaging over the finite family of spike realizations fixes a deterministic
instance with at least this expected regret for the given randomized
controller.
\end{proof}

In this construction, $H=2$ and $T=3n$. With $L_u>0$ fixed, the allowed
state marginals grow as large as $nL_u=(T/3)L_u$, while expected regret is
at least $L_uT/12$. Thus a bound depending only on $T$, $H$, and $L_u$
cannot be sublinear in $T$ uniformly over all allowed magnitudes of the
nonnegative state marginals.

\subsection{Minimax bounds and unresolved deadline dependence}
For each controller, consider its largest expected regret over admissible
arrival and cost sequences fixed before the run begins. The expectation
is over the controller's internal randomness. Taking the infimum of
these worst-case values over controllers that use only revealed
information and meet every arrival's deadline defines
$\mathfrak R_{T,H}(L_x,L_u)$. The admissible costs satisfy the joint-cost
assumptions with bounds $L_x$ and $L_u$.
For $H\ge3$, $T\ge4$, the established results imply
\begin{equation}\label{ch6:eq:minimax-sandwich}
 \begin{split}
 \max\left\{\frac{L_u}{2\sqrt3}\sqrt{T-1},\,
 \frac{L_x}{256\sqrt3}H^{3/2}\sqrt T\right\}
 &\le\mathfrak R_{T,H}(L_x,L_u),\\
 \mathfrak R_{T,H}(L_x,L_u)
 &\le\frac{2L_u+(H-1)L_x}{2}
 \sqrt{T-1}\,\Lambda_{H}.
 \end{split}
\end{equation}
The two subclasses give separate lower bounds, so their maximum bounds
the minimax regret from below.
In particular,
\begin{equation}\label{ch6:eq:state-minimax}
 \mathfrak R_{T,H}(L_x,0)
 =\Theta(L_xH^{3/2}\sqrt T)
 \quad(H\ge3,\ T\ge4).
\end{equation}
For fixed $L_x>0$ and bounded ratio $L_u/L_x$, the same dependence on
$H,T$ follows for the general class, with constants allowed to depend
on that ratio. These bounds do not establish matching dependence on
$L_x$ and $L_u$ when they vary independently.

For consumption-only costs, the separate conclusion is
\begin{equation}\label{ch6:eq:action-gap}
 \frac{L_u}{2\sqrt3}\sqrt{T-1}
 \le\mathfrak R_{T,H}(0,L_u)
 \le2(1+\sqrt2)L_u\sqrt{H(T-1)}
 \qquad(H\ge2,\ T\ge2).
\end{equation}
When costs depend only on consumption and the deadline is fixed, the
bounds establish the optimal $\sqrt T$ growth. Their dependence on the
deadline remains unresolved: the upper bound increases with
$\sqrt{H}$, while the lower bound has no deadline factor. The state
lower bound uses costs that depend on storage, so it cannot answer this
question for consumption-only costs. The sharp regret rate under the
additional nonnegative storage-marginal assumptions in
\eqref{ch6:eq:positive} also remains open.

\section{Future Research Directions}
\label{ch6:sec:next}
The main open question is the action contribution in
\eqref{ch6:eq:action-gap}. The current controller waits for an entire
interval before updating one interleaved copy, although a single observed
pair $(\zeta_t^x,\zeta_t^u)$ informs several active arrivals and intervals.

\begin{Thm:Remark}[Progressively revealed split feedback]
For a fixed $(s,I)$ with $I=[a,b]$, let $\pi^{\mathrm{L}}_i,\pi^{\mathrm{R}}_i$ be the final
conditional leaf fractions in its two children, extended by zero outside
the corresponding child, and define
\begin{equation}\label{ch6:eq:partial-coefficients}
 \delta_i=\pi^{\mathrm{L}}_i-\pi^{\mathrm{R}}_i,\qquad
 D_k=\sum_{i=k}^b\delta_i=-\sum_{i=a}^{k-1}\delta_i.
\end{equation}
The two expressions for $D_k$ agree because both conditional distributions
have unit mass and hence $\sum_i\delta_i=0$.
Expanding \eqref{ch6:eq:g} gives the exact identity
\begin{equation}\label{ch6:eq:partial}
 \boxed{g_{s,I}=w_s\left[
 \sum_{i=a}^b\delta_i\zeta_{s+i}^u
 +\sum_{k=a+1}^bD_k\zeta_{s+k}^x\right].}
\end{equation}
Each conditional leaf fraction is fixed before its service, and $D_k$ can
be computed from earlier conditional services through
\eqref{ch6:eq:partial-coefficients}. Hence the summands are revealed
progressively even though the full $g_{s,I}$ is unavailable at the split.
\end{Thm:Remark}

Using partial feedback requires a new regret analysis. Each revealed
piece of feedback belongs to a split that has already been chosen. An
early update changes the learner's reference value while some feedback
for that split is still missing. The analysis must therefore account for
the full cost comparison associated with the original split and count
each feedback contribution exactly once. Such an analysis might reduce
the cost of maintaining separate learners for delayed feedback.

Another direction is to prove a stronger lower bound for consumption-only
costs. A lower bound that grows with the deadline could show whether the
current $\sqrt{H}$ factor is necessary.

\section{Conclusion}
\label{ch6:sec:conclusion}

This paper studies online allocation under resource constraints and
individual service deadlines. The controller divides each arrival among
age intervals as they begin, ensuring complete consumption by the deadline
while allowing later decisions to use newly observed costs. Delayed
mirror-descent updates provide regret guarantees against the best fixed
age-allocation policy under adversarial jointly convex costs. Uniformly
positive state marginal costs can justify a shorter service deadline, and
known state costs can be incorporated directly into causal predictions.

A matching lower bound for signed linear storage costs shows that the
storage-dependent term in the general regret bound is unavoidable. The
optimal deadline dependence for consumption-only costs remains open.
Exploiting the progressively revealed feedback shared by active arrivals
and intervals is a direction for closing this gap.

\appendix
\section{Supporting Lemmas}

\subsection{Scalar quadratic and entropy geometry}
\label{app:geometry}

For a differentiable convex function $\psi$, its Bregman divergence is
\[
 D_\psi(p,z)=\psi(p)-\psi(z)-\psi'(z)(p-z).
\]
The quadratic regularizer $\psi_Q(p)=p^2/2$ gives
$D_{\psi_Q}(p,z)=(p-z)^2/2$ and
$\sup_{p\in[0,1]}D_{\psi_Q}(p,1/2)=1/8$.
For binary entropy, let
\[
 \psi_E(p)=p\log p+(1-p)\log(1-p),\qquad p\in[0,1],
\]
with the continuous convention at the endpoints. Its divergence is
$\operatorname{kl}(p\|z)$ for $z\in(0,1)$.
Since
\[
 \psi_E''(p)=\frac1p+\frac1{1-p}
 =\frac1{p(1-p)}\ge4\qquad(0<p<1),
\]
integration and continuity at the endpoints give
\[
 D_{\psi_E}(p,z)\ge2(p-z)^2,\qquad
 D_{\psi_E}(p,1/2)=\log2+\psi_E(p)\le\log2.
\]
These are the binary specializations of standard entropy geometry
bounds~\cite{BECK2003167}.

For the reference update, quadratic minimization gives
$z^+=[z-\eta g]_0^1$. Under entropy, first-order optimality gives
\[
 \eta g+\log\frac{z^+}{1-z^+}-\log\frac z{1-z}=0,
\]
hence
\[
 z^+=\frac{z e^{-\eta g}}{1-z+z e^{-\eta g}}.
\]
Strict convexity gives uniqueness. For finite $g$ and $z\in(0,1)$ this
update is interior. The prediction objective with a finite continuous
convex differentiable predictor also has an interior entropy minimizer:
the entropy derivative tends to $-\infty$ and $+\infty$ at the two
endpoints, respectively. Quadratic minimizers may lie on the boundary.

\subsection{Prediction and correction from a common reference}
\label{app:prediction}

The following scalar inequality is the prediction--correction comparison
used in the analysis; compare~\cite[Lemma~3]{RakhlinSridharan2013}.

\begin{Thm:Lemma}[Prediction and correction from a common reference]
\label{app:lem:optimistic}
Let $\psi$ be either regularizer in Appendix~\ref{app:geometry}, and let
$\sigma>0$ satisfy
$D_\psi(v,w)\ge \sigma(v-w)^2/2$ on its domain.
Let $\eta>0$, let $F:[0,1]\to\Real$ be convex and differentiable with
finite derivative, and let $g\in\Real$. Define
\[
 p=\argmin_{v\in[0,1]}\{\eta F(v)+D_\psi(v,z)\},\qquad
 z^+=\argmin_{v\in[0,1]}\{\eta gv+D_\psi(v,z)\}.
\]
For entropy assume $z\in(0,1)$; for the quadratic case allow $z\in[0,1]$.
Writing $\widehat g=F'(p)$, every $u\in[0,1]$ satisfies
\[
 g(p-u)\le
 \frac{D_\psi(u,z)-D_\psi(u,z^+)}{\eta}
 +\frac{\eta}{2\sigma}(g-\widehat g)^2
 -\frac{D_\psi(p,z)}{\eta}.
\]
\end{Thm:Lemma}
\begin{proof}
Direct expansion yields the three-point identity
\[
 (\psi'(a)-\psi'(b))(c-a)
 =D_\psi(c,b)-D_\psi(c,a)-D_\psi(a,b).
\]
The correction optimality condition, tested at $u$, therefore gives
\[
 \eta g(z^+-u)
 \le D_\psi(u,z)-D_\psi(u,z^+)-D_\psi(z^+,z).
\]
The prediction optimality condition, tested at $z^+$, similarly gives
\[
 \eta\widehat g(p-z^+)
 \le D_\psi(z^+,z)-D_\psi(z^+,p)-D_\psi(p,z).
\]
Both conditions allow quadratic boundary minimizers; the entropy
minimizers are interior by Appendix~\ref{app:geometry}.
Adding the inequalities and the remaining term
$\eta(g-\widehat g)(p-z^+)$ yields
\begin{align*}
 \eta g(p-u)
 &\le D_\psi(u,z)-D_\psi(u,z^+)-D_\psi(p,z)\\
 &\quad-D_\psi(z^+,p)+\eta(g-\widehat g)(p-z^+)\\
 &\le D_\psi(u,z)-D_\psi(u,z^+)-D_\psi(p,z)
       +\frac{\eta^2}{2\sigma}(g-\widehat g)^2.
\end{align*}
The last step uses strong convexity and
$\eta ab\le \sigma b^2/2+\eta^2a^2/(2\sigma)$.
Division by $\eta$ proves the claim.
\end{proof}

\subsection{Independent signs and deterministic instances}
\label{app:signs}

\begin{Thm:Lemma}[Weighted independent-sign bound; cf.~{\cite{haagerup1981khintchine}}]
\label{lem:intro:signs}\label{ch6:lem:signs}
For independent random signs $\xi_1,\ldots,\xi_n$, each equally likely to be
$-1$ or $+1$, and deterministic real coefficients $a_i$, let
$S=\sum_i a_i\xi_i$. Then
\begin{equation}\label{eq:intro:sign-moment}
 \E[|S|]\ge\frac1{\sqrt3}\left(\sum_i a_i^2\right)^{1/2}.
\end{equation}
In particular, when all $a_i=1$,
\begin{equation}\label{ch6:eq:signs}
 \E[|S|]\ge\sqrt{n/3}.
\end{equation}
\end{Thm:Lemma}
\begin{proof}
Put $v=\sum_i a_i^2$.  Independence and the fact that products containing an
odd power of any sign have zero expectation give
\begin{align*}
 \E[S^2]&=\sum_i a_i^2=v,\\
 \E[S^4]
 &=\sum_i a_i^4+6\sum_{i<j}a_i^2a_j^2\\
 &=3\left(\sum_i a_i^2\right)^2-2\sum_i a_i^4
 \le3v^2.
\end{align*}
H\"older's inequality with exponents $3/2$ and $3$ gives
\begin{align*}
 v=\E[|S|^2]
 &=\E\bigl[|S|^{2/3}|S|^{4/3}\bigr]\\
 &\le(\E[|S|])^{2/3}(\E[|S|^4])^{1/3}
 \le(\E[|S|])^{2/3}(3v^2)^{1/3}.
\end{align*}
For $v>0$, cubing and cancelling $v^2$ yields
$\E[|S|]\ge\sqrt{v/3}$; for $v=0$, both sides vanish.
\end{proof}

The classical Khintchine inequality gives the sharper bound
\cite{haagerup1981khintchine}:
\begin{equation}\label{eq:intro:khintchine}
 \E\left[\left|\sum_i a_i\xi_i\right|\right]
 \ge\frac1{\sqrt2}\left(\sum_i a_i^2\right)^{1/2}.
\end{equation}
To obtain a deterministic instance, first fix a controller, including the
rule for its random choices. Then draw an input sequence from a finite
collection, independently of those choices; each sequence specifies all
inputs before the run begins. If regret averaged over both the sequence
draw and the controller's random choices is at least $R$, then some
sequence in the collection has expected regret at least $R$ when only
the controller is randomized. This sequence can depend on which
controller was fixed.

\bibliographystyle{plain}
\bibliography{references}

\end{document}
